
\documentclass[aps,preprint]{revtex4}
%%%%%%%%%%%%%%%%%%%%%%%%%%%%%%%%%%%%%%%%%%%%%%%%%%%%%%%%%%%%%%%%%%%%%%%%%%%%%%%%%%%%%%%%%%%%%%%%%%%%%%%%%%%%%%%%%%%%%%%%%%%%
\usepackage{amsfonts}
\usepackage{amsmath}
\usepackage{amssymb}
\usepackage{graphicx}

\setcounter{MaxMatrixCols}{10}
%TCIDATA{OutputFilter=LATEX.DLL}
%TCIDATA{Version=4.10.0.2363}
%TCIDATA{Created=Friday, May 24, 2002 09:27:08}
%TCIDATA{LastRevised=Friday, September 12, 2003 14:13:51}
%TCIDATA{<META NAME="GraphicsSave" CONTENT="32">}
%TCIDATA{<META NAME="DocumentShell" CONTENT="Articles\SW\REVTeX 4">}
%TCIDATA{CSTFile=revtex4.cst}

\newtheorem{theorem}{Theorem}
\newtheorem{acknowledgement}[theorem]{Acknowledgement}
\newtheorem{algorithm}[theorem]{Algorithm}
\newtheorem{axiom}[theorem]{Axiom}
\newtheorem{claim}[theorem]{Claim}
\newtheorem{conclusion}[theorem]{Conclusion}
\newtheorem{condition}[theorem]{Condition}
\newtheorem{conjecture}[theorem]{Conjecture}
\newtheorem{corollary}[theorem]{Corollary}
\newtheorem{criterion}[theorem]{Criterion}
\newtheorem{definition}[theorem]{Definition}
\newtheorem{example}[theorem]{Example}
\newtheorem{exercise}[theorem]{Exercise}
\newtheorem{lemma}[theorem]{Lemma}
\newtheorem{notation}[theorem]{Notation}
\newtheorem{problem}[theorem]{Problem}
\newtheorem{proposition}[theorem]{Proposition}
\newtheorem{remark}[theorem]{Remark}
\newtheorem{solution}[theorem]{Solution}
\newtheorem{summary}[theorem]{Summary}
\newenvironment{proof}[1][Proof]{\noindent\textbf{#1.} }{\ \rule{0.5em}{0.5em}}
\input{tcilatex}

\begin{document}

\preprint{HamJac}
\title[HamJac]{Hamitlon-Jacobi Approach}
\author{Brian Weiner}
\affiliation{PSU}
\keywords{05/26}
\pacs{}

\begin{abstract}
Shell document for REV\TeX{} 4.
\end{abstract}

\startpage{1}
\endpage{1}
\maketitle
\tableofcontents

\section{ Details of Lagrangian and EOM}

\subsection{Parameterization of State Vectors}

$\mathcal{\breve{V}}\left( t\right) $The time dependence of a pure $N$ body
state can be described in terms of the non compact abelian semi group formed
by the elements 
\begin{equation}
\mathcal{\breve{V}}\left( t\right) e^{i\breve{\omega}\left( t\right)
}\equiv\int\mathcal{V}\left( \mathsf{z,}t\right) \breve{\Omega}_{N}\left( 
\mathsf{z}\right) d\mathsf{z}e^{i\int\omega\left( \mathsf{z,}t\right) \breve{%
\Omega}_{N}\left( \mathsf{z}\right) d\mathsf{z}}
\end{equation}
acting on a reference state $\left\vert \Psi_{ref}\right\rangle $ 
\begin{subequations}
\begin{align}
\left\vert \Psi\left( t\right) \right\rangle & =\mathcal{\breve{V}}\left(
t\right) e^{i\breve{\omega}\left( t\right) }\left\vert
\Psi_{ref}\right\rangle \\
\mathcal{\breve{V}}\left( t\right) & =\int\mathcal{V}\left( \mathsf{z,}%
t\right) \breve{\Omega}_{N}\left( \mathsf{z}\right) d\mathsf{z} \\
\breve{\omega}\left( t\right) & =\int\omega\left( \mathsf{z,}t\right) \breve{%
\Omega}_{N}\left( \mathsf{z}\right) d\mathsf{z} \\
\Psi\left( \mathsf{z,}t\right) & =\left\langle \mathsf{z}|\Psi\left(
t\right) \right\rangle =\left\langle \mathsf{z}|\mathcal{\breve{V}}\left(
t\right) e^{i\breve{\omega}\left( t\right) }\Psi_{ref}\right\rangle \\
\Psi\left( \mathsf{z,}t\right) & =\Psi\left( \mathcal{\varkappa}\left( 
\mathsf{z,}t\right) \right) =\Psi\left( \mathcal{V}\left( \mathsf{z,}%
t\right) ,\omega\left( \mathsf{z,}t\right) \right) =\mathcal{V}\left( 
\mathsf{z,}t\right) e^{i\omega\left( \mathsf{z,}t\right) }\Psi_{ref}\left( 
\mathsf{z}\right)
\end{align}
where $\mathcal{\varkappa}\left( \mathsf{z,}t\right) \equiv\left( \mathcal{V}%
\left( \mathsf{z,}t\right) ,\omega\left( \mathsf{z,}t\right) \right) $.

\subsection{Density Operators}

The density operators associated with these states are 
\end{subequations}
\begin{equation}
D_{\Psi}^{N}\left( t\right) =\frac{\left\vert \Psi\left( t\right)
\right\rangle \left\langle \Psi\left( t\right) \right\vert }{\left\langle
\Psi\left( t\right) |\Psi\left( t\right) \right\rangle }
\end{equation}
this trajectory of density operators is clearly invariant to complex scaling
of the vector $\left\vert \Psi\left( t\right) \right\rangle $ i.e. 
\begin{equation}
D_{\Psi}^{N}\left( t\right) =\frac{\left\vert \eta\left( t\right) \Psi\left(
t\right) \right\rangle \left\langle \eta\left( t\right) \Psi\left( t\right)
\right\vert }{\left\langle \eta\left( t\right) \Psi\left( t\right)
|\eta\left( t\right) \Psi\left( t\right) \right\rangle }\qquad\forall\eta%
\left( t\right) \in\mathbb{C}
\end{equation}
where $\eta\left( t\right) $ is a purely time dependent $\lceil$i.e. $\frac{%
\partial\eta\left( t\right) }{_{\partial\mathsf{z}}}=0;\frac {%
\partial\eta\left( t\right) }{_{\partial t}}\neq0\rfloor$ complex valued
function. This of course is consistent with the postulates of quantum
mechanics where the pure states correspond to lines (i.e. one dimensional
subspaces) in a complex Hilbert space. A pure state then corresponds to an
equivalence class of vectors that is formed by the vectors that belong to a
given line. Any vector belonging to that line is a representative of the
equivalence class and can be said to be a \emph{representative pure state
vector} which we will shorten to \emph{state vector }where the words
\textquotedblright representative\textquotedblright\ and \textquotedblright
pure\textquotedblright\ are to be implicitly understood.\emph{\ }The initial
and final pure state density operators of a quantum process are only
uniquely defined in terms of lines belonging to a complex Hilbert space i.e.
lines generated by initial and final state vectors. Hence a necessary
component of the action principle that determines quantum trajectories in
terms of \emph{state vectors} corresponds to a calculus of variation problem
with initial and final state vectors belonging to lines in a complex Hilbert
space. The solution of such problems is found by finding stationary
processes for the fixed endpoint problem between a fixed initial and a fixed
final point on these lines then fixing the position on these lines by
applying \emph{tranversality constraints}. However in order to obtain the
Time Dependent Schr\"{o}dinger Equation (TDSE) it turns out this is not a
sufficient condition and one must consider stationary process that take
different lengths of time while proceeding between the initial and final
state vectors i.e. the time duration of the process is allowed to vary and
one must select trajectories that are stationary for the fixed end point
problem, lie on the given initial and final lines belonging to $\mathcal{H}%
^{N}$ \emph{and are also stationary with respect to the duration of the
process.}

The action is defined as%
\begin{equation}
\mathcal{A}\left( \Psi ,\tau \right) =\int_{t_{i}}^{t_{i}+\tau }\mathcal{L}%
\left( \Psi ,\dot{\Psi},t\right) dt
\end{equation}%
where 
\begin{equation}
\mathcal{L}\left( \Psi ,\dot{\Psi},t\right) =\frac{\left\langle \Psi \left(
t\right) |\left( i\frac{\partial }{\partial t}-H\right) \Psi \left( t\right)
\right\rangle }{\left\langle \Psi \left( t\right) |\Psi \left( t\right)
\right\rangle }
\end{equation}%
The possible boundary conditions on the action are come from considering:

\begin{description}
\item[$\bullet $] What are the properties of the paths $\left\vert \Psi
\right\rangle :\left[ t_{i},t_{i}+\tau \right] \equiv \left[ t_{i},t_{f}%
\right] \rightarrow \mathcal{H}^{N}$ i.e. are they continuous, bounded ...?

\item[$\bullet $] Are the end points $\left\vert \Psi \left( t_{i}\right)
\right\rangle =\left\vert \Phi _{1}\right\rangle ;\left\vert \Psi \left(
t_{f}\right) \right\rangle =\left\vert \Phi _{2}\right\rangle $ fixed,
constrained to lie in a given manifold of $\mathcal{H}^{N}$ or unconstrained?

\item[$\bullet $] Is the time interval $\tau $ fixed, variable or
constrained in any manner?
\end{description}

We resolve these questions in the following manner:

\begin{description}
\item[$\bullet $] In standard quantum mechanics the paths should be strongly
continuous i.e. for each $t\in \left[ t_{i},t_{i}+\tau \right] \equiv \bar{%
\Upsilon}$ and for every $\epsilon $ there exists a $\delta $ such that 
\begin{equation}
\left\vert t^{\prime }-t\right\vert <\delta \Rightarrow \left\Vert
\left\vert \Psi \left( t^{\prime }\right) \right\rangle -\left\vert \Psi
\left( t\right) \right\rangle \right\Vert <\epsilon
\end{equation}%
for all $t^{\prime }\in \left[ t_{i},t_{i}+\tau \right] $, which also
implies 
\begin{equation}
s-\lim_{t_{k}\rightarrow t}\left\vert \Psi \left( t_{k}\right) \right\rangle
=\left\vert \Psi \left( t\right) \right\rangle
\end{equation}%
for all sequences $\left\{ t_{k}\right\} $ converging to $t$.

\item[$\bullet $] The initial point is fixed i.e. 
\begin{equation}
\left\vert \Psi \left( t_{i}\right) \right\rangle =\left\vert \Phi
_{1}\right\rangle
\end{equation}%
while the end point is constrained to lie on a fixed ray in $\mathcal{H}^{N}$
i.e.%
\begin{equation}
\left\vert \Psi \left( t_{f}\right) \right\rangle =\eta \left\vert \Phi
_{2}\right\rangle ;\quad \Phi _{2}\text{ is fixed, }\eta =\rho e^{i\sigma
}\in \mathbb{C}
\end{equation}

\item[$\bullet $] The time interval is unconstrained, but finite i.e. $\bar{%
\Upsilon}$ is compact.

\item The arguments of the action can be partitioned into the independent
variables 
\begin{eqnarray}
\left\vert \Psi \left( t\right) \right\rangle \text{ for }t &\in &\Upsilon
\equiv \left( t_{i},t_{i}+\tau \right) ;\,\left\vert \Psi \left(
t_{i}\right) \right\rangle \equiv \left\vert \Phi _{1}\right\rangle ;  \notag
\\
\,\left\vert \Psi \left( t_{i}+\tau \right) \right\rangle &\equiv &\eta
\left\vert \Phi _{2}\right\rangle \text{ and }\tau  \label{part}
\end{eqnarray}%
and one obtains that 
\begin{equation}
\mathcal{A}\left( \Psi ,\tau \right) =\mathcal{A}\left( \left. \Psi
\right\vert _{\Upsilon },\Phi _{1},\Phi _{2},\eta ,\tau \right)
\end{equation}%
$\lceil $If we want the paths $\left\vert \Psi \left( t\right) \right\rangle 
$ to be strongly continuous then we must consider spaces of paths, $\mathcal{%
P}\left( \left\vert \Phi _{1}\right\rangle ,\eta \left\vert \Phi
_{2}\right\rangle \right) $, that are strongly continuous that have end
point values $\left\vert \Phi _{1}\right\rangle $ and $\eta \left\vert \Phi
_{2}\right\rangle $. For a fixed $t\in \Upsilon $ every point in $\mathcal{H}
$ lies on a path $\left\vert \Psi \left( t\right) \right\rangle $ belonging
to a fixed $\mathcal{P}\left( \left\vert \Phi _{1}\right\rangle ,\eta
\left\vert \Phi _{2}\right\rangle \right) $ thus varying over the values, $%
\left\vert \delta \Psi \left( t\right) \right\rangle $, of every strongly
continuous increment path, $\left\vert \delta \Psi \right\rangle $ at a
fixed $t$ is the same as varying over all vectors in $\mathcal{H}$. At the
end points the only allowable variaotions are 
\begin{eqnarray}
\left\vert \delta \Psi \left( t_{i}\right) \right\rangle &=&0  \notag \\
\left\vert \delta \Psi \left( t_{f}\right) \right\rangle &=&\eta ^{\prime
}\left\vert \Phi _{2}\right\rangle
\end{eqnarray}%
$\rfloor $
\end{description}

Eq. $\left( \ref{part}\right) $ leads to the following expression for the
first order variation%
\begin{align}
\delta ^{\left( 1\right) }\mathcal{A}& =\int \frac{\delta \mathcal{A}\left(
\Psi ,\tau \right) }{\delta \left. \Psi \right\vert _{\Upsilon }\left( 
\mathsf{z}\right) }\delta \Psi \left( t^{\prime },\mathsf{z}\right) d\mathsf{%
z}dt^{\prime }+\int \frac{\delta \mathcal{A}\left( \Psi ,\tau \right) }{%
\delta \Psi _{i}\left( \mathsf{z}\right) }\delta \Psi _{i}\left( \mathsf{z}%
\right) d\mathsf{z}  \notag \\
& \qquad \qquad \qquad +\int \frac{\delta \mathcal{A}\left( \Psi ,\tau
\right) }{\delta \Psi _{f}\left( \mathsf{z}\right) }\delta \Psi _{f}\left( 
\mathsf{z}\right) d\mathsf{z}+\frac{\partial \mathcal{A}\left( \Psi ,\tau
\right) }{\partial \tau }\delta \tau  \label{DA}
\end{align}%
The condition for \emph{complete }stationary action is then 
\begin{equation}
\frac{\delta \mathcal{A}\left( \Psi ,\tau \right) }{\delta \left. \Psi
\right\vert _{\Upsilon }\left( \mathsf{z}\right) }=0;\frac{\delta \mathcal{A}%
\left( \Psi ,\tau \right) }{\delta \Psi _{i}\left( \mathsf{z}\right) }=0;%
\frac{\delta \mathcal{A}\left( \Psi ,\tau \right) }{\delta \Psi _{f}\left( 
\mathsf{z}\right) }=0;\frac{\partial \mathcal{A}\left( \Psi ,\tau \right) }{%
\partial t}=0
\end{equation}%
$\lceil $Note that 
\begin{equation}
\frac{\partial \mathcal{A}}{\partial \tau }=\frac{\partial \mathcal{A}}{%
\partial t}
\end{equation}%
$\rfloor $. The first two conditions can be replaced by $\frac{\delta 
\mathcal{\dot{A}}\left( \Psi \left( t\right) \right) }{\delta \Psi \left( 
\mathsf{z}\right) }=0$ as described below.

One can also observe that the total rate of change of the complete action is 
\begin{equation}
\frac{d\mathcal{A}}{dt}=\int \frac{\delta \mathcal{A}\left( \Psi ,\tau
\right) }{\delta \left. \Psi \right\vert _{\Upsilon }\left( \mathsf{z}%
\right) }\delta \Psi \left( t,\mathsf{z}\right) d\mathsf{z}+\frac{\partial 
\mathcal{A}\left( \Psi ,\tau \right) }{\partial t}
\end{equation}

\subsection{Variational Derivative}

In terms of the Lagragian the complete first order variation is 
\begin{equation}
\delta ^{\left( 1\right) }\mathcal{A}=\int_{t_{i}}^{t_{i}+\tau }\left\{ 
\frac{\delta \mathcal{L}\left( \Psi ,\dot{\Psi},t\right) }{\delta \Psi
\left( \mathsf{z}\right) }\delta \Psi \left( t,\mathsf{z}\right) +\frac{%
\delta \mathcal{L}\left( \Psi ,\dot{\Psi},t\right) }{\delta \dot{\Psi}\left( 
\mathsf{z}\right) }\delta \dot{\Psi}\left( t,\mathsf{z}\right) \right\} d%
\mathsf{z}dt+\frac{\partial \mathcal{A}\left( \Psi ,\tau \right) }{\partial
\tau }\delta \tau
\end{equation}%
giving 
\begin{align}
\delta ^{\left( 1\right) }\mathcal{A}& =\int_{t_{i}}^{t_{i}+\tau }\left\{ 
\frac{\delta \mathcal{L}\left( \Psi ,\dot{\Psi},\tau \right) }{\delta \Psi
\left( \mathsf{z}\right) }-\frac{d}{dt}\frac{\delta \mathcal{L}\left( \Psi ,%
\dot{\Psi},\tau \right) }{\delta \dot{\Psi}\left( \mathsf{z}\right) }%
\right\} \delta \Psi \left( \mathsf{z}\right) d\mathsf{z}dt  \notag \\
& +\frac{\delta \mathcal{L}\left( \Psi ,\dot{\Psi},t_{i}+\tau \right) }{%
\delta \dot{\Psi}\left( \mathsf{z}\right) }\delta \Psi _{f}\left( \mathsf{z}%
\right) -\frac{\delta \mathcal{L}\left( \Psi ,\dot{\Psi},t_{i}\right) }{%
\delta \dot{\Psi}\left( \mathsf{z}\right) }\delta \Psi _{i}\left( \mathsf{z}%
\right)  \notag \\
& \qquad \qquad \qquad \qquad \qquad \qquad \qquad \qquad \qquad +\frac{%
\partial \mathcal{A}\left( \Psi ,\tau \right) }{\partial \tau }\delta \tau
\end{align}%
and hence 
\begin{align}
\delta ^{\left( 1\right) }\mathcal{A}& =\int_{t_{i}}^{t_{i}+\tau }\frac{%
\delta \mathcal{\dot{A}}\left( \Psi \left( t\right) \right) }{\delta \Psi
\left( \mathsf{z}\right) }\delta \Psi \left( \mathsf{z}\right) d\mathsf{z}dt
\notag \\
& +\frac{\delta \mathcal{L}\left( \Psi ,\dot{\Psi},t_{i}+\tau \right) }{%
\delta \dot{\Psi}\left( \mathsf{z}\right) }\delta \Psi _{f}\left( \mathsf{z}%
\right) -\frac{\delta \mathcal{L}\left( \Psi ,\dot{\Psi},t_{i}\right) }{%
\delta \dot{\Psi}\left( \mathsf{z}\right) }\delta \Psi _{i}\left( \mathsf{z}%
\right)  \notag \\
& \qquad \qquad \qquad \qquad \qquad \qquad \qquad \qquad \qquad +\frac{%
\partial \mathcal{A}\left( \Psi ,\tau \right) }{\partial \tau }\delta \tau
\end{align}%
where the \textit{variational derivative }is defined as 
\begin{equation}
\frac{\delta \mathcal{\dot{A}}\left( \Psi \left( t\right) \right) }{\delta
\Psi \left( \mathsf{z}\right) }=\frac{\delta \mathcal{L}\left( \Psi ,\dot{%
\Psi},t\right) }{\delta \Psi \left( \mathsf{z}\right) }-\frac{d}{dt}\left\{ 
\frac{\delta \mathcal{L}\left( \Psi ,\dot{\Psi},t\right) }{\delta \dot{\Psi}%
\left( \mathsf{z}\right) }\right\}
\end{equation}

\subsection{Fixed End Point Problem}

We first then consider the fixed end point problem. A stationary time
evolution of a state vector is determined by the fixed end point equation of
motion obtained when the variational derivative $\frac{\delta \mathcal{\dot{A%
}}\left( \Psi \left( t\right) \right) }{\delta \Psi \left( \mathsf{z}\right) 
}$ is set equal to zero 
\begin{equation}
\frac{\delta \mathcal{\dot{A}}\left( \Psi \left( t\right) \right) }{\delta
\Psi \left( \mathsf{z}\right) }=\frac{\delta \mathcal{L}\left( \Psi ,\dot{%
\Psi},t\right) }{\delta \Psi \left( \mathsf{z}\right) }-\frac{d}{dt}\left\{ 
\frac{\delta \mathcal{L}\left( \Psi ,\dot{\Psi},t\right) }{\delta \dot{\Psi}%
\left( \mathsf{z}\right) }\right\} =0  \label{eom1}
\end{equation}%
which is obtained from the stationarity condition plus the boundary
condition of fixed end points 
\begin{align}
\delta ^{\left( 1\right) }\mathcal{A}\left( \Psi ,\tau \right) &
=\int_{t_{i}}^{t_{f}}\delta ^{\left( 1\right) }\mathcal{\dot{A}}\left( \Psi
\left( t\right) \right) dt=\int_{t_{i}}^{t_{f}}\delta ^{\left( 1\right) }%
\mathcal{L}\left( \Psi ,\dot{\Psi},t\right) dt=0  \notag \\
& \left. {}\right. \qquad \qquad \qquad \qquad \forall \delta \Psi \left(
t\right) ;\quad \delta \Psi \left( t_{f}\right) =\delta \Psi \left(
t_{i}\right) =0
\end{align}%
for the action%
\begin{equation}
\mathcal{A}\left( \Psi ,\tau \right) =\int_{t_{i}}^{t_{f}}\mathcal{\dot{A}}%
\left( \Psi \left( t\right) \right) dt=\int_{t_{i}}^{t_{f}}\mathcal{L}\left(
\Psi ,\dot{\Psi},t\right) dt
\end{equation}%
where 
\begin{equation}
\delta ^{\left( 1\right) }\mathcal{\dot{A}}\left( \Psi \left( t\right)
\right) =\int \frac{\delta \mathcal{\dot{A}}\left( \Psi \left( t\right)
\right) }{\delta \Psi \left( \mathsf{z}\right) }\delta \Psi \left( \mathsf{z,%
}t\right) d\mathsf{z}=\left\langle \left. \left[ \frac{\delta \mathcal{\dot{A%
}}\left( \Psi \left( t\right) \right) }{\delta \Psi }\right] ^{\ast
}\right\vert \delta \Psi \left( t\right) \right\rangle
\end{equation}%
$\lceil $In group parameterized form eq. $\left( \ref{eom1}\right) $ is 
\begin{equation}
\frac{\delta \mathcal{\dot{A}}\left( \varkappa \left( t\right) \right) }{%
\delta \varkappa \left( \mathsf{z}\right) }=\frac{\delta \mathcal{L}\left( 
\mathcal{\varkappa },\mathcal{\dot{\varkappa}},t\right) }{\delta \mathcal{%
\varkappa }\left( \mathsf{z}\right) }-\frac{d}{dt}\left\{ \frac{\delta 
\mathcal{L}\left( \mathcal{\varkappa },\mathcal{\dot{\varkappa}},t\right) }{%
\delta \mathcal{\dot{\varkappa}}\left( \mathsf{z}\right) }\right\} =0
\label{eom2}
\end{equation}%
$\rfloor $

\subsection{Lagrangian}

We consider a Lagrangian $\mathcal{L}$ given by 
\begin{align}
\mathcal{L}\left( \Psi\left( t\right) ,\dot{\Psi}\left( t\right) \right) & =i%
\frac{\left\langle \Psi\left( t\right) |\dot{\Psi}\left( t\right)
\right\rangle }{\left\langle \Psi\left( t\right) |\Psi\left( t\right)
\right\rangle }-\frac{\left\langle \Psi\left( t\right) |H\Psi\left( t\right)
\right\rangle }{\left\langle \Psi\left( t\right) |\Psi\left( t\right)
\right\rangle }  \notag \\
& =\frac{\left\langle \Psi\left( t\right) |\left( i\frac{\partial}{\partial t%
}-H\right) \Psi\left( t\right) \right\rangle }{\left\langle \Psi\left(
t\right) |\Psi\left( t\right) \right\rangle }  \notag \\
& =\frac{\int\left\{ i\Psi^{\ast}\left( \mathsf{z,}t\right) \dot{\Psi }%
\left( \mathsf{z,}t\right) -\Psi^{\ast}\left( \mathsf{z,}t\right) \left[
H\Psi\right] \left( \mathsf{z,}t\right) \right\} }{\left\langle \Psi\left(
t\right) |\Psi\left( t\right) \right\rangle }d\mathsf{z}  \label{Lar}
\end{align}
Note that there is a many to one relationship between Lagrangians and a EOM.
A class (but not all) of Lagrangians that produce the same equation of
motion can be obtained from 
\begin{equation}
\mathcal{L}_{F}\left( \Psi\left( t\right) ,\dot{\Psi}\left( t\right) \right)
=\mathcal{L}\left( \Psi\left( t\right) ,\dot{\Psi}\left( t\right) \right) +%
\frac{dF\left( \Psi\left( t\right) ,\dot{\Psi}\left( t\right) \right) }{dt}
\label{LF}
\end{equation}
as then 
\begin{align}
& \mathcal{A}_{F}\left( \Psi,\tau\right) =  \notag \\
& \qquad\qquad\qquad\int_{t_{i}}^{t_{f}}\mathcal{L}_{F}\left( \Psi\left(
t\right) ,\dot{\Psi}\left( t\right) \right) dt=\int_{t_{i}}^{t_{f}}\mathcal{L%
}\left( \Psi\left( t\right) ,\dot{\Psi}\left( t\right) \right) dt  \notag \\
& \qquad\qquad\qquad\qquad\qquad\qquad\qquad\qquad\qquad+F\left( \Psi\left(
t_{f}\right) ,\dot{\Psi}\left( t_{f}\right) \right) -F\left( \Psi\left(
t_{i}\right) ,\dot{\Psi}\left( t_{i}\right) \right)
\end{align}
and 
\begin{equation}
\delta^{\left( 1\right) }\mathcal{A}_{F}\left( \Psi,\tau\right)
=\int_{t_{i}}^{t_{f}}\delta^{\left( 1\right) }\mathcal{L}\left( \Psi ,\dot{%
\Psi},t\right) dt=\int_{t_{i}}^{t_{f}}\delta^{\left( 1\right) }\mathcal{\dot{%
A}}\left( \Psi\left( t\right) \right) dt=\delta^{\left( 1\right) }\mathcal{A}%
\left( \Psi,\tau\right)
\end{equation}
since $F\left( \Psi\left( t_{f}\right) ,\dot{\Psi}\left( t_{f}\right)
\right) $ and $F\left( \Psi\left( t_{i}\right) ,\dot{\Psi}\left(
t_{i}\right) \right) $ are constant as the end points are fixed.

We can define $F$ as 
\begin{equation}
F\left( \Psi\left( t\right) ,\dot{\Psi}\left( t\right) \right) =-\frac{i}{2}%
\ln\left\{ \frac{\left\langle \Psi\left( t\right) |\Psi\left( t\right)
\right\rangle }{\left\langle \Psi\left( t_{i}\right) |\Psi\left(
t_{i}\right) \right\rangle }\right\} =-\frac{i}{2}\int_{t_{i}}^{t}\frac {%
d\ln\left\langle \Psi\left( t^{\prime}\right) |\Psi\left( t^{\prime }\right)
\right\rangle }{dt^{\prime}}dt^{\prime}
\end{equation}
and thus obtain the modified action 
\begin{equation}
\mathcal{A}_{F}\left( \Psi,\tau\right) =\int_{t_{i}}^{t_{f}}\mathcal{L}%
_{F}\left( \Psi\left( t\right) ,\dot{\Psi}\left( t\right) \right) dt=%
\mathcal{A}\left( \Psi,\tau\right) -\frac{i}{2}\ln\left\{ \frac {%
\left\langle \Psi\left( t_{f}\right) |\Psi\left( t_{f}\right) \right\rangle 
}{\left\langle \Psi\left( t_{i}\right) |\Psi\left( t_{i}\right)
\right\rangle }\right\}  \label{AF}
\end{equation}
where 
\begin{align}
& \mathcal{L}_{F}\left( \Psi\left( t\right) ,\dot{\Psi}\left( t\right)
\right) =\mathcal{L}\left( \Psi\left( t\right) ,\dot{\Psi}\left( t\right)
\right) -\frac{i}{2}\frac{d\ln\left\langle \Psi\left( t\right) |\Psi\left(
t\right) \right\rangle }{dt}  \notag \\
& \left. {}\right. =i\frac{\left\langle \Psi\left( t\right) |\dot{\Psi }%
\left( t\right) \right\rangle }{\left\langle \Psi\left( t\right) |\Psi\left(
t\right) \right\rangle }-\frac{\left\langle \Psi\left( t\right) |H\Psi\left(
t\right) \right\rangle }{\left\langle \Psi\left( t\right) |\Psi\left(
t\right) \right\rangle }-\frac{i}{2}\left\{ \frac{\left\langle \Psi\left(
t\right) |\dot{\Psi}\left( t\right) \right\rangle +\left\langle \dot{\Psi}%
\left( t\right) |\Psi\left( t\right) \right\rangle }{\left\langle \Psi\left(
t\right) |\Psi\left( t\right) \right\rangle }\right\}  \notag \\
& \left. {}\right. \qquad\qquad\qquad\qquad=\frac{i}{2}\frac{\left\{
\left\langle \Psi\left( t\right) |\dot{\Psi}\left( t\right) \right\rangle
-\left\langle \dot{\Psi}\left( t\right) |\Psi\left( t\right) \right\rangle
\right\} }{\left\langle \Psi\left( t\right) |\Psi\left( t\right)
\right\rangle }-\frac{\left\langle \Psi\left( t\right) |H\Psi\left( t\right)
\right\rangle }{\left\langle \Psi\left( t\right) |\Psi\left( t\right)
\right\rangle }  \label{LF1}
\end{align}
One can further develop eq. $\left( \ref{LF1}\right) $ to obtain 
\begin{equation}
\mathcal{L}_{F}\left( \Psi\left( t\right) ,\dot{\Psi}\left( t\right) \right)
=-\left\{ \frac{\func{Im}\left\langle \Psi\left( t\right) |\dot{\Psi}\left(
t\right) \right\rangle }{\left\langle \Psi\left( t\right) |\Psi\left(
t\right) \right\rangle }+\frac{\left\langle \Psi\left( t\right) |H\Psi\left(
t\right) \right\rangle }{\left\langle \Psi\left( t\right) |\Psi\left(
t\right) \right\rangle }\right\} =\func{Re}\mathcal{L}\left( \Psi\left(
t\right) ,\dot{\Psi}\left( t\right) \right)
\end{equation}
and thus observe that 
\begin{equation}
\mathcal{A}_{F}\left( \Psi,\tau\right) =\int_{t_{i}}^{t_{f}}\func{Re}%
\mathcal{L}\left( \Psi\left( t\right) ,\dot{\Psi}\left( t\right) \right) dt=%
\func{Re}\mathcal{A}\left( \Psi,\tau\right)
\end{equation}
and deduce that 
\begin{equation}
\func{Im}\mathcal{A}\left( \Psi,\tau\right) =\frac{1}{2}\ln\left\{ \frac{%
\left\langle \Psi\left( t_{f}\right) |\Psi\left( t_{f}\right) \right\rangle 
}{\left\langle \Psi\left( t_{i}\right) |\Psi\left( t_{i}\right)
\right\rangle }\right\}
\end{equation}
Obviously 
\begin{equation}
\delta^{\left( 1\right) }\mathcal{A}_{F}\left( \Psi,\tau\right)
=0\Longleftrightarrow\delta^{\left( 1\right) }\mathcal{A}\left( \Psi
,\tau\right) =0
\end{equation}
since the extra terms in eq. $\left( \ref{AF}\right) $ are constant as the
paths have fixed end points. Thus $\left| \Psi_{s}\right\rangle $ is a
stationary point of $\mathcal{A}\left( \Psi,\tau\right) $ iff it is a
stationary point of $\mathcal{A}_{F}\left( \Psi,\tau\right) $ if the
constraint of fixed endpoints is invoked. We can summarize these
relationships by 
\begin{align}
\mathcal{L}\left( \Psi\left( t\right) ,\dot{\Psi}\left( t\right) \right) & =%
\func{Re}\mathcal{L}\left( \Psi\left( t\right) ,\dot{\Psi }\left( t\right)
\right) +i\func{Im}\mathcal{L}\left( \Psi\left( t\right) ,\dot{\Psi}\left(
t\right) \right)  \notag \\
& =-\frac{1}{\mathcal{N}\left( \Psi\left( t\right) \right) }\left\{ \func{Im}%
\left\langle \Psi\left( t\right) |\dot{\Psi}\left( t\right) \right\rangle
+\left\langle \Psi\left( t\right) |H\Psi\left( t\right) \right\rangle
\right\} +\frac{i}{2}\frac{d\ln\left\langle \Psi\left( t\right) |\Psi\left(
t\right) \right\rangle }{dt}  \notag \\
& =\mathcal{L}_{F}\left( \Psi\left( t\right) ,\dot{\Psi}\left( t\right)
\right) +\frac{i}{2}\frac{d\ln\left\langle \Psi\left( t\right) |\Psi\left(
t\right) \right\rangle }{dt}
\end{align}
and 
\begin{align}
\mathcal{A}\left( \Psi,\tau\right) & =\func{Re}\mathcal{A}\left(
\Psi,\tau\right) +i\func{Im}\mathcal{A}\left( \Psi,\tau\right)
\qquad\qquad\qquad\qquad\qquad\qquad\qquad\qquad\qquad\qquad\qquad  \notag \\
& =\int_{t_{i}}^{t_{f}}\func{Re}\mathcal{L}\left( \Psi\left( t\right) ,\dot{%
\Psi}\left( t\right) \right) dt+\frac{i}{2}\ln\left\{ \frac{\left\langle
\Psi\left( t_{f}\right) |\Psi\left( t_{f}\right) \right\rangle }{%
\left\langle \Psi\left( t_{i}\right) |\Psi\left( t_{i}\right) \right\rangle }%
\right\}  \notag \\
& =\mathcal{A}_{F}\left( \Psi,\tau\right) +\frac{i}{2}\ln\left\{ \frac{%
\left\langle \Psi\left( t_{f}\right) |\Psi\left( t_{f}\right) \right\rangle 
}{\left\langle \Psi\left( t_{i}\right) |\Psi\left( t_{i}\right)
\right\rangle }\right\}
\end{align}

\subsection{Computation of Variational Derivatives}

In order to calculate the variational derivative of $\mathcal{A}\left(
\Psi,\tau\right) $ we first tabulate the following definitions and
derivatives 
\begin{align}
\mathcal{T}\left( \Psi\left( t\right) ,\dot{\Psi}\left( t\right) \right) &
=i\left\langle \Psi\left( t\right) |\dot{\Psi}\left( t\right) \right\rangle 
\notag \\
\mathcal{T}_{F}\left( \Psi\left( t\right) ,\dot{\Psi}\left( t\right) \right)
& =\frac{i}{2}\left\{ \left\langle \Psi\left( t\right) |\dot {\Psi}\left(
t\right) \right\rangle -\left\langle \dot{\Psi}\left( t\right) |\Psi\left(
t\right) \right\rangle \right\}  \notag \\
\mathcal{H}\left( \Psi\left( t\right) \right) & =\left\langle \Psi\left(
t\right) |H\Psi\left( t\right) \right\rangle  \notag \\
\mathcal{N}\left( \Psi\left( t\right) \right) & =\left\langle \Psi\left(
t\right) |\Psi\left( t\right) \right\rangle  \notag \\
\mathcal{L}\left( \Psi\left( t\right) ,\dot{\Psi}\left( t\right) \right) & =%
\frac{\mathcal{T}\left( \Psi\left( t\right) ,\dot{\Psi}\left( t\right)
\right) -\mathcal{H}\left( \Psi\left( t\right) \right) }{\mathcal{N}\left(
\Psi\left( t\right) \right) }  \notag \\
\mathcal{L}_{F}\left( \Psi\left( t\right) ,\dot{\Psi}\left( t\right) \right)
& =\frac{\mathcal{T}_{F}\left( \Psi\left( t\right) ,\dot{\Psi }\left(
t\right) \right) -\mathcal{H}\left( \Psi\left( t\right) \right) }{\mathcal{N}%
\left( \Psi\left( t\right) \right) }
\end{align}
which give the following derivatives 
\begin{align}
\frac{\delta\mathcal{N}\left( \Psi\left( t\right) \right) }{\delta
\Psi\left( \mathsf{z}\right) } & =\Psi^{\ast}\left( \mathsf{z,}t\right) 
\notag \\
& \equiv\frac{\delta\mathcal{N}\left( \Psi\left( t\right) \right) }{%
\delta\left| \Psi\right\rangle }=\left\langle \Psi\left( t\right) \right| \\
\frac{\delta\mathcal{T}_{F}\left( \Psi\left( t\right) ,\dot{\Psi}\left(
t\right) \right) }{\delta\Psi\left( \mathsf{z}\right) } & =-\frac{i}{2}\dot{%
\Psi}^{\ast}\left( \mathsf{z,}t\right)  \notag \\
& \equiv\frac{\delta\mathcal{T}_{F}\left( \Psi\left( t\right) ,\dot{\Psi }%
\left( t\right) \right) }{\delta\left| \Psi\right\rangle }=-\frac{i}{2}%
\left\langle \dot{\Psi}\left( t\right) \right| \\
\frac{\delta\mathcal{T}_{F}\left( \Psi\left( t\right) ,\dot{\Psi}\left(
t\right) \right) }{\delta\dot{\Psi}\left( \mathsf{z}\right) } & =\frac {i}{2}%
\Psi^{\ast}\left( \mathsf{z,}t\right)  \notag \\
& \equiv\frac{i}{2}\left\langle \Psi\left( t\right) \right| \\
\frac{\delta\mathcal{H}\left( \Psi\left( t\right) \right) }{\delta
\Psi\left( \mathsf{z}\right) } & =\int\Psi^{\ast}\left( \mathsf{z}^{\prime}%
\mathsf{,}t\right) H\left( \mathsf{z}^{\prime},\mathsf{z}\right) d\mathsf{z}%
^{\prime}\equiv\left\langle H\Psi\left( t\right) |\mathsf{z}\right\rangle =%
\left[ H\Psi\right] ^{\ast}\left( \mathsf{z,}t\right)  \notag \\
& \equiv\frac{\delta\mathcal{H}\left( \Psi\left( t\right) \right) }{%
\delta\left| \Psi\right\rangle }=\left\langle H\Psi\left( t\right) \right|
=\left\langle \Psi\left( t\right) \right| H \\
\frac{\delta\mathcal{T}\left( \Psi\left( t\right) ,\dot{\Psi}\left( t\right)
\right) }{\delta\Psi\left( \mathsf{z}\right) } & =0  \notag \\
& \equiv\frac{\delta\mathcal{T}\left( \Psi\left( t\right) ,\dot{\Psi }\left(
t\right) \right) }{\delta\left| \Psi\right\rangle } \\
\frac{\delta\mathcal{T}\left( \Psi\left( t\right) ,\dot{\Psi}\left( t\right)
\right) }{\delta\dot{\Psi}\left( \mathsf{z}\right) } & =i\Psi^{\ast}\left( 
\mathsf{z,}t\right)  \notag \\
& \equiv i\left\langle \Psi\left( t\right) \right|
\end{align}
Initially consider the first term of eq.$\left( \ref{eom1}\right) $ using
the Lagrangian eq. $\left( \ref{Lar}\right) $ and obtain 
\begin{align}
& \frac{\delta\mathcal{L}\left( \Psi\left( t\right) ,\dot{\Psi}\left(
t\right) \right) }{\delta\Psi\left( \mathsf{z}\right) }=  \notag \\
& \qquad\frac{1}{\mathcal{N}\left( \Psi\left( t\right) \right) }\left\{ 
\frac{\delta\left\{ \mathcal{T}\left( \Psi\left( t\right) ,\dot{\Psi }\left(
t\right) \right) \mathcal{-H}\left( \Psi\left( t\right) \right) \right\} }{%
\delta\Psi\left( \mathsf{z}\right) }-\mathcal{L}\left( \Psi\left( t\right) ,%
\dot{\Psi}\left( t\right) \right) \frac {\delta\mathcal{N}\left( \Psi\left(
t\right) \right) }{\delta\Psi\left( \mathsf{z}\right) }\right\}
\end{align}
where%
\begin{align}
\frac{\delta\mathcal{L}_{F}\left( \Psi\left( t\right) ,\dot{\Psi}\left(
t\right) \right) }{\delta\dot{\Psi}\left( \mathsf{z}\right) } & =\frac{%
i\Psi^{\ast}\left( \mathsf{z,}t\right) }{2\mathcal{N}\left( \Psi\left(
t\right) \right) }  \notag \\
& \equiv\frac{\delta\mathcal{L}_{F}\left( \Psi\left( t\right) ,\dot{\Psi }%
\left( t\right) \right) }{\delta\left| \dot{\Psi}\right\rangle }=\frac {i}{2%
\mathcal{N}\left( \Psi\left( t\right) \right) }\left\langle \Psi\left(
t\right) \right| \\
\frac{\delta\mathcal{L}\left( \Psi\left( t\right) ,\dot{\Psi}\left( t\right)
\right) }{\delta\dot{\Psi}\left( \mathsf{z}\right) } & =i\frac{%
\Psi^{\ast}\left( \mathsf{z,}t\right) }{\mathcal{N}\left( \Psi\left(
t\right) \right) }  \notag \\
& \equiv\frac{\delta\mathcal{L}\left( \Psi\left( t\right) ,\dot{\Psi }\left(
t\right) \right) }{\delta\left| \dot{\Psi}\right\rangle }=\frac {i}{\mathcal{%
N}\left( \Psi\left( t\right) \right) }\left\langle \Psi\left( t\right)
\right| \\
\frac{\delta\mathcal{L}_{F}\left( \Psi\left( t\right) ,\dot{\Psi}\left(
t\right) \right) }{\delta\Psi\left( \mathsf{z}\right) } & =-\frac {1}{%
\mathcal{N}\left( \Psi\left( t\right) \right) }\left\{ \frac{i}{2}\dot{\Psi}%
\left( \mathsf{z,}t\right) ^{\ast}+\left[ H\Psi\right] ^{\ast}\left( \mathsf{%
z,}t\right) +\mathcal{L}_{F}\left( \Psi\left( t\right) ,\dot{\Psi}\left(
t\right) \right) \Psi^{\ast}\left( \mathsf{z,}t\right) \right\}  \notag \\
\left. \equiv\right. \frac{\delta\mathcal{L}_{F}\left( \Psi\left( t\right) ,%
\dot{\Psi}\left( t\right) \right) }{\delta\left| \Psi \right\rangle } & =-%
\frac{1}{\mathcal{N}\left( \Psi\left( t\right) \right) }\left\{ \frac{i}{2}%
\left\langle \dot{\Psi}\right| +\left\langle \Psi\right| H+\mathcal{L}%
_{F}\left( \Psi\left( t\right) ,\dot{\Psi}\left( t\right) \right)
\left\langle \Psi\right| \right\} \\
\frac{\delta\mathcal{L}\left( \Psi\left( t\right) ,\dot{\Psi}\left( t\right)
\right) }{\delta\Psi\left( \mathsf{z}\right) } & =-\frac {1}{\mathcal{N}%
\left( \Psi\left( t\right) \right) }\left\{ \left[ H\Psi\right]
^{\ast}\left( \mathsf{z,}t\right) +\mathcal{L}\left( \Psi\left( t\right) ,%
\dot{\Psi}\left( t\right) \right) \Psi^{\ast}\left( \mathsf{z,}t\right)
\right\}  \notag \\
& \equiv\frac{\delta\mathcal{L}\left( \Psi\left( t\right) ,\dot{\Psi }\left(
t\right) \right) }{\delta\left| \Psi\right\rangle }=-\frac {1}{\mathcal{N}%
\left( \Psi\left( t\right) \right) }\left\langle \Psi\right| \left\{ H+%
\mathcal{L}\left( \Psi\left( t\right) ,\dot{\Psi }\left( t\right) \right)
I\right\}
\end{align}
In order to calculate $\frac{\delta\mathcal{\dot{A}}\left( \Psi\left(
t\right) \right) }{\delta\Psi\left( \mathsf{z}\right) }$ we need only to
observe that 
\begin{equation}
\frac{d}{dt}\left\{ \frac{\delta\mathcal{L}\left( \Psi,\dot{\Psi},t\right) }{%
\delta\dot{\Psi}\left( \mathsf{z}\right) }\right\} =\frac{i}{\mathcal{N}%
\left( \Psi\left( t\right) \right) }\left\{ \dot{\Psi}^{\ast }\left( \mathsf{%
z,}t\right) -\frac{\Psi^{\ast}\left( \mathsf{z,}t\right) }{\mathcal{N}\left(
\Psi\left( t\right) \right) }\frac{d\mathcal{N}\left( \Psi\left( t\right)
\right) }{dt}\right\}
\end{equation}
which is \ equivalent to 
\begin{equation}
\frac{d}{dt}\left\{ \frac{\delta\mathcal{L}\left( \Psi,\dot{\Psi},t\right) }{%
\delta\left| \dot{\Psi}\right\rangle }\right\} =\frac{i}{\mathcal{N}\left(
\Psi\left( t\right) \right) }\left\{ \left\langle \dot{\Psi }\right| -\frac{%
\left\langle \Psi\right| }{\mathcal{N}\left( \Psi\left( t\right) \right) }%
\frac{d\mathcal{N}\left( \Psi\left( t\right) \right) }{dt}\right\}
\end{equation}
which produces 
\begin{align}
& \frac{\delta\mathcal{\dot{A}}\left( \Psi\left( t\right) \right) }{%
\delta\Psi\left( \mathsf{z}\right) }=  \notag \\
& \frac{1}{\mathcal{N}\left( \Psi\left( t\right) \right) }\left\{ \frac{%
\delta\mathcal{T}\left( \Psi\left( t\right) ,\dot{\Psi}\left( t\right)
\right) }{\delta\Psi\left( \mathsf{z}\right) }-\left[ H\Psi\right]
^{\ast}\left( \mathsf{z,}t\right) -\mathcal{L}\left( \Psi\left( t\right) ,%
\dot{\Psi}\left( t\right) \right) \Psi^{\ast}\left( \mathsf{z,}t\right)
\right\}  \notag \\
& \qquad\qquad\qquad\qquad\qquad\qquad\qquad\qquad\qquad\qquad\qquad \qquad-%
\frac{d}{dt}\left\{ \frac{\delta\mathcal{L}\left( \Psi\left( t\right) ,\dot{%
\Psi}\left( t\right) \right) }{\delta\dot{\Psi}\left( \mathsf{z}\right) }%
\right\}  \notag \\
&  \notag \\
& =-\frac{1}{\mathcal{N}\left( \Psi\left( t\right) \right) }\left\{ \left[
H\Psi\right] ^{\ast}\left( \mathsf{z,}t\right) +\mathcal{L}\left( \Psi\left(
t\right) ,\dot{\Psi}\left( t\right) \right) \Psi^{\ast}\left( \mathsf{z,}%
t\right) \right\}  \notag \\
& \qquad\qquad\qquad\qquad\qquad\qquad\qquad-\frac{i}{\mathcal{N}\left(
\Psi\left( t\right) \right) }\left\{ \dot{\Psi}^{\ast}\left( \mathsf{z,}%
t\right) -\frac{\Psi^{\ast}\left( \mathsf{z,}t\right) }{\mathcal{N}\left(
\Psi\left( t\right) \right) }\frac{d\mathcal{N}\left( \Psi\left( t\right)
\right) }{dt}\right\}  \notag \\
&  \notag \\
& =-\frac{1}{\mathcal{N}\left( \Psi\left( t\right) \right) }\left\{ i\dot{%
\Psi}^{\ast}\left( \mathsf{z,}t\right) \right. +\left[ H\Psi\right]
^{\ast}\left( \mathsf{z,}t\right)  \notag \\
& \qquad\left. +\left\{ \frac{i\left\langle \Psi\left( t\right) |\dot {\Psi}%
\left( t\right) \right\rangle -\left\langle \Psi\left( t\right) |H\Psi\left(
t\right) \right\rangle -i\left\langle \Psi\left( t\right) |\dot{\Psi}\left(
t\right) \right\rangle -i\left\langle \dot{\Psi}\left( t\right) |\Psi\left(
t\right) \right\rangle }{\mathcal{N}\left( \Psi\left( t\right) \right) }%
\right\} \Psi^{\ast}\left( \mathsf{z,}t\right) \right\}  \notag \\
&  \notag \\
& =-\frac{1}{\mathcal{N}\left( \Psi\left( t\right) \right) }\left\{ i\dot{%
\Psi}^{\ast}\left( \mathsf{z,}t\right) +\left[ H\Psi\right] ^{\ast }\left( 
\mathsf{z,}t\right) -\frac{\left\{ i\left\langle \dot{\Psi}\left( t\right)
|\Psi\left( t\right) \right\rangle +\left\langle \Psi\left( t\right)
|H\Psi\left( t\right) \right\rangle \Psi^{\ast}\left( \mathsf{z,}t\right)
\right\} }{\mathcal{N}\left( \Psi\left( t\right) \right) }\right\}  \notag \\
& =-\frac{1}{\mathcal{N}\left( \Psi\left( t\right) \right) }\left\{ i\frac{%
\partial\Psi^{\ast}\left( \mathsf{z,}t\right) }{\partial t}+\int H\left( 
\mathsf{z,z}^{\prime}\right) ^{\ast}\Psi^{\ast}\left( \mathsf{z}^{\prime}%
\mathsf{,}t\right) d\mathsf{z}^{\prime}-\mathcal{L}\left( \Psi\left(
t\right) ,\dot{\Psi}\left( t\right) \right) ^{\ast}\Psi^{\ast }\left( 
\mathsf{z,}t\right) \right\}
\end{align}
which is equivalent to 
\begin{equation}
\frame{$\frac{\delta\mathcal{\dot{A}}\left( \Psi\left( t\right) \right) }{%
\delta\left| \Psi\right\rangle }$=-$\frac{1}{\mathcal{N}\left( \Psi\left(
t\right) \right) }\left\{ i\left\langle \dot{\Psi}\right| +\left\langle
H\Psi\right| +\mathcal{L}\left( \Psi\left( t\right) ,\dot{\Psi}\left(
t\right) \right) ^\ast\left\langle \Psi\right| \right\} $}
\end{equation}
In order to evaluate the variational derivative $\frac{\delta\mathcal{\dot{A}%
}_{F}\left( \Psi\left( t\right) \right) }{\delta\Psi\left( \mathsf{z}\right) 
} $ we first consider 
\begin{align}
& \frac{\delta\mathcal{L}_{F}\left( \Psi\left( t\right) ,\dot{\Psi}\left(
t\right) \right) }{\delta\Psi\left( \mathsf{z}\right) }  \notag \\
& \qquad\left. =\right. -\frac{1}{\mathcal{N}\left( \Psi\left( t\right)
\right) }\left\{ \frac{i}{2}\dot{\Psi}^{\ast}\left( \mathsf{z,}t\right) +%
\left[ H\Psi\right] ^{\ast}\left( \mathsf{z,}t\right) +\mathcal{L}_{F}\left(
\Psi\left( t\right) ,\dot{\Psi}\left( t\right) \right) \Psi^{\ast}\left( 
\mathsf{z,}t\right) \right\}  \label{a} \\
& \equiv\frac{\delta\mathcal{L}_{F}\left( \Psi\left( t\right) ,\dot{\Psi }%
\left( t\right) \right) }{\delta\left| \Psi\right\rangle }  \notag \\
& \qquad\left. =\right. -\frac{1}{\mathcal{N}\left( \Psi\left( t\right)
\right) }\left\{ \frac{i}{2}\left\langle \dot{\Psi}\left( t\right) \right|
+\left\langle H\Psi\left( t\right) \right| +\mathcal{L}_{F}\left( \Psi\left(
t\right) ,\dot{\Psi}\left( t\right) \right) \left\langle \Psi\left( t\right)
\right| \right\}
\end{align}
The we consider the second term of eq.$\left( \ref{eom1}\right) $ and
observe that 
\begin{align}
& \frac{d}{dt}\left\{ \frac{\delta\mathcal{L}_{F}\left( \Psi\left( t\right) ,%
\dot{\Psi}\left( t\right) \right) }{\delta\dot{\Psi}\left( \mathsf{z}\right) 
}\right\} =\frac{i}{2}\frac{1}{\mathcal{N}\left( \Psi\left( t\right) \right) 
}\left\{ \dot{\Psi}^{\ast}\left( \mathsf{z,}t\right) -\frac{%
\Psi^{\ast}\left( \mathsf{z,}t\right) }{\mathcal{N}\left( \Psi\left(
t\right) \right) }\frac{d\mathcal{N}\left( \Psi\left( t\right) \right) }{dt}%
\right\}  \label{b} \\
& \equiv\frac{d}{dt}\left\{ \frac{\delta\mathcal{L}_{F}\left( \Psi\left(
t\right) ,\dot{\Psi}\left( t\right) \right) }{\delta\left| \dot{\Psi }%
\right\rangle }\right\} =\frac{i}{2}\frac{1}{\mathcal{N}\left( \Psi\left(
t\right) \right) }\left\{ \left\langle \Psi\left( t\right) \right| -\frac{%
\left\langle \Psi\left( t\right) \right| }{\mathcal{N}\left( \Psi\left(
t\right) \right) }\frac{d\mathcal{N}\left( \Psi\left( t\right) \right) }{dt}%
\right\}
\end{align}
combining eqs. $\left( \ref{a}\right) $ and $\left( \ref{b}\right) $ we
produce 
\begin{align}
& \frac{\delta\mathcal{\dot{A}}_{F}\left( \Psi\left( t\right) \right) }{%
\delta\Psi\left( \mathsf{z}\right) }=  \notag \\
& -\frac{1}{\mathcal{N}\left( \Psi\left( t\right) \right) }\left\{ \frac{%
\delta\mathcal{T}_{F}\left( \Psi\left( t\right) ,\dot{\Psi}\left( t\right)
\right) }{\delta\Psi\left( \mathsf{z}\right) }+\left[ H\Psi\right]
^{\ast}\left( \mathsf{z,}t\right) +\mathcal{L}_{F}\left( \Psi\left( t\right)
,\dot{\Psi}\left( t\right) \right) \Psi^{\ast}\left( \mathsf{z,}t\right)
\right\}  \notag \\
& \qquad\qquad\qquad\qquad\qquad\qquad\qquad\qquad\qquad\qquad\qquad-\frac {d%
}{dt}\left\{ \frac{\delta\mathcal{L}_{F}\left( \Psi\left( t\right) ,\dot{\Psi%
}\left( t\right) \right) }{\delta\dot{\Psi}\left( \mathsf{z}\right) }\right\}
\notag \\
& \left. =\right. -\frac{1}{\mathcal{N}\left( \Psi\left( t\right) \right) }%
\left\{ \frac{i}{2}\dot{\Psi}^{\ast}\left( \mathsf{z,}t\right) +\left[ H\Psi%
\right] ^{\ast}\left( \mathsf{z,}t\right) +\mathcal{L}_{F}\left( \Psi\left(
t\right) ,\dot{\Psi}\left( t\right) \right) \Psi^{\ast}\left( \mathsf{z,}%
t\right) \right\}  \notag \\
& \qquad\qquad\qquad\qquad\qquad\qquad-\frac{i}{2}\frac{1}{\mathcal{N}\left(
\Psi\left( t\right) \right) }\left\{ \dot{\Psi}^{\ast}\left( \mathsf{z,}%
t\right) -\frac{\Psi^{\ast}\left( \mathsf{z,}t\right) }{\mathcal{N}\left(
\Psi\left( t\right) \right) }\frac{d\mathcal{N}\left( \Psi\left( t\right)
\right) }{dt}\right\}
\end{align}
After noting that 
\begin{align}
& \frac{\mathcal{L}_{F}\left( \Psi\left( t\right) ,\dot{\Psi}\left( t\right)
\right) \Psi^{\ast}\left( \mathsf{z,}t\right) }{\mathcal{N}\left( \Psi\left(
t\right) \right) }  \notag \\
& \left. =\right. \frac{i}{2}\frac{\left\{ \left\langle \Psi\left( t\right) |%
\dot{\Psi}\left( t\right) \right\rangle -\left\langle \dot{\Psi }\left(
t\right) |\Psi\left( t\right) \right\rangle \right\} \Psi^{\ast }\left( 
\mathsf{z,}t\right) }{\mathcal{N}\left( \Psi\left( t\right) \right) ^{2}}-%
\frac{\left\langle \Psi\left( t\right) |H\Psi\left( t\right) \right\rangle
\Psi^{\ast}\left( \mathsf{z,}t\right) }{\mathcal{N}\left( \Psi\left(
t\right) \right) ^{2}}
\end{align}
and 
\begin{equation}
\frac{i}{2}\frac{\Psi^{\ast}\left( \mathsf{z,}t\right) }{\mathcal{N}\left(
\Psi\left( t\right) \right) }\frac{d\mathcal{N}\left( \Psi\left( t\right)
\right) }{dt}=\frac{i}{2}\frac{\left\{ \left\langle \Psi\left( t\right) |%
\dot{\Psi}\left( t\right) \right\rangle +\left\langle \dot{\Psi }\left(
t\right) |\Psi\left( t\right) \right\rangle \right\} \Psi^{\ast }\left( 
\mathsf{z,}t\right) }{\mathcal{N}\left( \Psi\left( t\right) \right) ^{2}}
\end{equation}
we obtain that 
\begin{align}
& -\frac{\mathcal{L}_{F}\left( \Psi\left( t\right) ,\dot{\Psi}\left(
t\right) \right) \Psi^{\ast}\left( \mathsf{z,}t\right) }{\mathcal{N}\left(
\Psi\left( t\right) \right) }+\frac{i}{2}\frac{\Psi^{\ast}\left( \mathsf{z,}%
t\right) }{\mathcal{N}\left( \Psi\left( t\right) \right) ^{2}}\frac{d%
\mathcal{N}\left( \Psi\left( t\right) \right) }{dt}  \notag \\
& \qquad\qquad\left. =\right. \frac{1}{\mathcal{N}\left( \Psi\left( t\right)
\right) }\left\{ \frac{i\left\langle \dot{\Psi}\left( t\right) |\Psi\left(
t\right) \right\rangle }{\mathcal{N}\left( \Psi\left( t\right) \right) }+%
\frac{\left\langle \Psi\left( t\right) |H\Psi\left( t\right) \right\rangle }{%
\mathcal{N}\left( \Psi\left( t\right) \right) }\right\} \Psi^{\ast}\left( 
\mathsf{z,}t\right)  \notag \\
& \qquad\qquad\qquad\qquad\qquad\quad\qquad\qquad\qquad\qquad\left. =\right.
-\frac{\mathcal{L}\left( \Psi\left( t\right) ,\dot{\Psi}\left( t\right)
\right) ^{\ast}\Psi^{\ast}\left( \mathsf{z,}t\right) }{\mathcal{N}\left(
\Psi\left( t\right) \right) }
\end{align}
and collecting all terms together we get the following expression for the
variational derivatives of $\mathcal{A}_{F}$ 
\begin{equation}
\frac{\delta\mathcal{\dot{A}}_{F}\left( \Psi\left( t\right) \right) }{%
\delta\Psi\left( \mathsf{z}\right) }=-\frac{1}{\mathcal{N}\left( \Psi\left(
t\right) \right) }\left\{ i\dot{\Psi}^{\ast}\left( \mathsf{z,}t\right) +%
\left[ H\Psi\right] ^{\ast}\left( \mathsf{z,}t\right) +\mathcal{L}\left(
\Psi\left( t\right) ,\dot{\Psi}\left( t\right) \right)
^{\ast}\Psi^{\ast}\left( \mathsf{z,}t\right) \right\}
\end{equation}
which is equivalent to 
\begin{equation}
\frame{$\frac{\delta\mathcal{\dot{A}}_F\left( \Psi\left( t\right) \right) }{%
\delta\left| \Psi\right\rangle }$=-$\frac{1}{\mathcal{N}\left( \Psi\left(
t\right) \right) }\left\{ i\left\langle \dot{\Psi}\right| +\left\langle
H\Psi\right| +\mathcal{L}\left( \Psi\left( t\right) ,\dot{\Psi}\left(
t\right) \right) ^\ast\left\langle \Psi\right| \right\} $}
\end{equation}
We have thus shown explicitly that 
\begin{equation}
\frac{\delta\mathcal{\dot{A}}\left( \Psi\left( t\right) \right) }{%
\delta\Psi\left( \mathsf{z}\right) }=\frac{\delta\mathcal{\dot{A}}_{F}\left(
\Psi\left( t\right) \right) }{\delta\Psi\left( \mathsf{z}\right) }\equiv%
\frac{\delta\mathcal{\dot{A}}\left( \Psi\left( t\right) \right) }{%
\delta\left| \Psi\right\rangle }=\frac{\delta\mathcal{\dot{A}}_{F}\left(
\Psi\left( t\right) \right) }{\delta\left| \Psi\right\rangle }.
\end{equation}

\subsection{Stationarity Condition}

Stationary points of the action wrt variation of the trajectories \emph{and}
the boundary conditions of fixed time duration, $\tau$, and endpoints $%
\left\{ \left| \Lambda_{i}\right\rangle ,\left| \Lambda_{f}\right\rangle
\right\} $ occur at trajectories where the variational derivative is zero
i.e. 
\begin{align}
& \frac{\delta\mathcal{\dot{A}}\left( \Psi\left( t\right) \right) }{%
\delta\Psi\left( \mathsf{z}\right) }=\frac{\delta\mathcal{\dot{A}}_{F}\left(
\Psi\left( t\right) \right) }{\delta\Psi\left( \mathsf{z}\right) }=0  \notag
\\
& \left. \Rightarrow\right. -i\dot{\Psi}^{\ast}\left( \mathsf{z,}t\right) -%
\left[ H\Psi\right] ^{\ast}\left( \mathsf{z,}t\right) =\mathcal{L}\left(
\Psi\left( t\right) ,\dot{\Psi}\left( t\right) \right) ^{\ast}\Psi^{\ast
}\left( \mathsf{z,}t\right)  \notag \\
& \left. \equiv\right. i\frac{\partial\Psi^{\ast}\left( \mathsf{z,}t\right) 
}{\partial t}+\int H\left( \mathsf{z,z}^{\prime}\right) ^{\ast
}\Psi^{\ast}\left( \mathsf{z}^{\prime}\mathsf{,}t\right) d\mathsf{z}%
^{\prime}=-\mathcal{L}\left( \Psi\left( t\right) ,\dot{\Psi}\left( t\right)
\right) ^{\ast}\Psi^{\ast}\left( \mathsf{z,}t\right)  \notag \\
& \left. \equiv\right. i\frac{\partial\Psi\left( \mathsf{z,}t\right) }{%
\partial t}-\int H\left( \mathsf{z,z}^{\prime}\right) \Psi\left( \mathsf{z}%
^{\prime}\mathsf{,}t\right) d\mathsf{z}^{\prime}=\mathcal{L}\left(
\Psi\left( t\right) ,\dot{\Psi}\left( t\right) \right) \Psi\left( \mathsf{z,}%
t\right)  \label{preSEden}
\end{align}
which is equivalent to 
\begin{align}
& \left\langle \left( i\frac{\partial}{\partial t}-H\right) \Psi\left(
t\right) \right| =\left\langle \mathcal{L}\left( \Psi\left( t\right) ,\dot{%
\Psi}\left( t\right) \right) \Psi\left( t\right) \right|  \notag \\
& \Rightarrow\left( i\frac{\partial}{\partial t}-H\right) \left| \Psi\left(
t\right) \right\rangle =\mathcal{L}\left( \Psi\left( t\right) ,\dot{\Psi}%
\left( t\right) \right) \left| \Psi\left( t\right) \right\rangle  \notag \\
& \Rightarrow i\frac{\partial}{\partial t}\left| \Psi\left( t\right)
\right\rangle =\left( H+\mathcal{L}\left( \Psi\left( t\right) ,\dot{\Psi }%
\left( t\right) \right) \right) \left| \Psi\left( t\right) \right\rangle
\label{preSE}
\end{align}
We can write explicitly write down the solution to eq. $\left( \ref{preSE}%
\right) $ 
\begin{align}
\left| \Psi\left( t\right) \right\rangle & =e^{-i\left\{ Ht\,+\,\int
_{t_{i}}^{t}\mathcal{L}\left( \Psi\left( t^{\prime}\right) ,\dot{\Psi }%
\left( t^{\prime}\right) \right) dt^{\prime}\right\} }\left| \Psi\left(
t_{i}\right) \right\rangle  \notag \\
& =e^{-i\left\{ Ht\,+\,\mathcal{A}\left( \Psi,\left[ t_{i},t\right] \right)
\right\} }\left| \Psi\left( t_{i}\right) \right\rangle  \notag \\
& =e^{-i\left\{ Ht\,+\,\mathcal{A}_{F}\left( \Psi,\left[ t_{i},t\right]
\right) +\frac{i}{2}\ln\left\{ \frac{\left\langle \Psi\left( t\right)
|\Psi\left( t\right) \right\rangle }{\left\langle \Psi\left( t_{i}\right)
|\Psi\left( t_{i}\right) \right\rangle }\right\} \right\} }\left| \Psi\left(
t_{i}\right) \right\rangle  \label{solution}
\end{align}
Note that $\left| \Psi\left( t_{i}\right) \right\rangle =\left|
\Lambda_{i}\right\rangle $ and $\left| \Psi\left( t\right) \right\rangle
=\left| \Lambda_{t}\right\rangle $ are arbitrary but fixed vectors.

\subsection{The Schr\"{o}dinger Equation}

The stationary action conditions $\delta\mathcal{A}\left( \Psi,\tau\right)
=0 $ for the boundary conditions of fixed time duration, $\tau$, and
endpoints $\left\{ \left| \Lambda_{i}\right\rangle ,\left|
\Lambda_{f}\right\rangle \right\} $ or even for the boundary conditions of
fixed time duration, $\tau $, and variable endpoints belonging to the lines $%
\left\{ \eta\left( t_{i}\right) \left| \Lambda_{i}\right\rangle ,\eta\left(
t_{f}\right) \left| \Lambda_{f}\right\rangle \right\} $ do not produce
solutions to the TDSE. In this section we display how to obtain solutions
from the from the fixed end point solutions by complex scaling of the
trajectories. If we define 
\begin{align}
\Phi\left( \mathsf{z,}t\right) & =\Psi\left( \mathsf{z,}t\right)
e^{i\int_{t_{i}}^{t}\mathcal{L}\left( \Psi\left( t^{\prime}\right) ,\dot{\Psi%
}\left( t^{\prime}\right) \right) dt^{\prime}}  \notag \\
& =\Psi\left( \mathsf{z,}t\right) e^{i\int_{t_{i}}^{t}\mathcal{L}\left(
\Psi\left( t^{\prime}\right) ,\dot{\Psi}\left( t^{\prime}\right) \right)
dt^{\prime}}  \notag \\
& =\Psi\left( \mathsf{z,}t\right) e^{i\int_{t_{i}}^{t}\left\{ \func{Re}%
\mathcal{L}\left( \Psi\left( t^{\prime}\right) ,\dot{\Psi }\left(
t^{\prime}\right) \right) +i\func{Im}\mathcal{L}\left( \Psi\left(
t^{\prime}\right) ,\dot{\Psi}\left( t^{\prime}\right) \right) \right\}
dt^{\prime}}  \notag \\
& =\Psi\left( \mathsf{z,}t\right) e^{i\int_{t_{i}}^{t}\func{Re}\mathcal{L}%
\left( \Psi\left( t^{\prime}\right) ,\dot{\Psi}\left( t^{\prime}\right)
\right) dt^{\prime}}e^{-\int_{t_{i}}^{t}\func{Im}\mathcal{L}\left(
\Psi\left( t^{\prime}\right) ,\dot{\Psi}\left( t^{\prime}\right) \right)
dt^{\prime}}  \notag \\
& =\Psi\left( \mathsf{z,}t\right) e^{i\int_{t_{i}}^{t}\func{Re}\mathcal{L}%
\left( \Psi\left( t^{\prime}\right) ,\dot{\Psi}\left( t^{\prime}\right)
\right) dt^{\prime}}e^{\ln\left\{ \sqrt{\frac {\left\langle \Psi\left(
t_{i}\right) |\Psi\left( t_{i}\right) \right\rangle }{\left\langle
\Psi\left( t\right) |\Psi\left( t\right) \right\rangle }}\right\} }  \notag
\\
& =\Psi\left( \mathsf{z,}t\right) e^{i\mathcal{A}\left( \Psi,\left[ t_{i},t%
\right] \right) }
\end{align}
then 
\begin{equation}
\dot{\Phi}\left( \mathsf{z,}t\right) =\dot{\Psi}\left( \mathsf{z,}t\right)
e^{i\int_{t_{i}}^{t}\mathcal{L}\left( \Psi,\dot{\Psi},t^{\prime}\right)
dt^{\prime}}+i\Psi\left( \mathsf{z,}t\right) \mathcal{L}\left( \Psi ,\dot{%
\Psi},t\right) e^{i\int_{t_{i}}^{t}\mathcal{L}\left( \Psi,\dot{\Psi }%
,t^{\prime}\right) dt^{\prime}}
\end{equation}
leading to%
\begin{align}
& =\dot{\Psi}\left( \mathsf{z,}t\right) e^{i\int_{t_{i}}^{t}\mathcal{L}%
\left( \Psi\left( t^{\prime}\right) ,\dot{\Psi}\left( t^{\prime}\right)
\right) dt^{\prime}}  \notag \\
& +i\left\{ i\frac{\partial\Psi\left( \mathsf{z,}t\right) }{\partial t}-\int
H\left( \mathsf{z,z}^{\prime}\right) \Psi\left( \mathsf{z}^{\prime }\mathsf{,%
}t\right) d\mathsf{z}^{\prime}\right\} e^{i\int_{t_{i}}^{t}\mathcal{L}\left(
\Psi\left( t^{\prime}\right) ,\dot{\Psi}\left( t^{\prime}\right) \right)
dt^{\prime}}  \notag \\
& =-i\left[ H\Psi\right] \left( \mathsf{z,}t\right) e^{-i\int_{t_{i}}^{t}%
\mathcal{L}\left( \Psi\left( t^{\prime}\right) ,\dot{\Psi}\left(
t^{\prime}\right) \right) dt^{\prime}}  \notag \\
& \Rightarrow\dot{\Phi}\left( \mathsf{z,}t\right) =-iH\Phi\left( \mathsf{z,}%
t\right) \equiv i\frac{\partial\Phi\left( \mathsf{z,}t\right) }{\partial t}%
=H\Phi\left( \mathsf{z,}t\right) \equiv\left( i\frac {\partial\Phi\left( 
\mathsf{z,}t\right) }{\partial t}-H\Phi\left( \mathsf{z,}t\right) \right) =0
\notag \\
& \equiv\left\{ i\frac{\partial}{\partial t}-H\right\} \left| \Phi\left(
t\right) \right\rangle =0  \label{SE}
\end{align}
which is the Time Dependent Schr\"{o}dinger Equation.

One can also note that eq.$\left( \ref{SE}\right) $ implies that 
\begin{equation}
\mathcal{L}\left( \Phi\left( t\right) ,\dot{\Phi}\left( t\right) \right) =0%
\text{ }
\end{equation}
$\lceil$In fact $\mathcal{L}\left( \Phi\left( t\right) ,\dot{\Phi}\left(
t\right) \right) =0$ plus stationarity for paths of fixed end points implies
the TDSE$\rfloor$ and that the norm of $\left| \Phi\left( t\right)
\right\rangle $ is constant as 
\begin{align}
\int\Phi^{\ast}\left( \mathsf{z,}t\right) \Phi\left( \mathsf{z,}t\right) d%
\mathsf{z} & \mathsf{=}\int\Psi^{\ast}\left( \mathsf{z,}t\right) \Psi\left( 
\mathsf{z,}t\right) e^{i\int_{t_{i}}^{t}\left\{ \mathcal{L}\left( \Psi\left(
t^{\prime}\right) ,\dot{\Psi}\left( t^{\prime}\right) \right) -\mathcal{L}%
\left( \Psi\left( t^{\prime}\right) ,\dot{\Psi}\left( t^{\prime}\right)
\right) ^{\ast}\right\} di}d\mathsf{z}  \notag \\
& =\int\Psi^{\ast}\left( \mathsf{z,}t\right) \Psi\left( \mathsf{z,}t\right)
e^{-\int_{t_{i}}^{t}\frac{d\ln\left\langle \Psi\left( t\right) |\Psi\left(
t\right) \right\rangle }{dt}di}d\mathsf{z}  \notag \\
& =\int\Psi^{\ast}\left( \mathsf{z,}t\right) \Psi\left( \mathsf{z,}t\right)
e^{-\ln\left\{ \frac{\left\langle \Psi\left( t\right) |\Psi\left( t\right)
\right\rangle }{\left\langle \Psi\left( t_{i}\right) |\Psi\left(
t_{i}\right) \right\rangle }\right\} }d\mathsf{z}  \notag \\
& =\int\Psi^{\ast}\left( \mathsf{z,}t\right) \Psi\left( \mathsf{z,}t\right) d%
\mathsf{z}\frac{\left\langle \Psi\left( t_{i}\right) |\Psi\left(
t_{i}\right) \right\rangle }{\left\langle \Psi\left( t\right) |\Psi\left(
t\right) \right\rangle }  \notag \\
& =\left\langle \Psi\left( t_{i}\right) |\Psi\left( t_{i}\right)
\right\rangle
\end{align}
However clearly 
\begin{equation}
\Phi\left( \mathsf{z,}t_{f}\right) \neq\Psi\left( \mathsf{z,}t_{f}\right)
\equiv\left| \Phi\left( t_{f}\right) \right\rangle \neq\left| \Psi\left(
t_{f}\right) \right\rangle
\end{equation}
since 
\begin{align}
& \Phi\left( \mathsf{z,}t_{f}\right) =\Psi\left( \mathsf{z,}t_{f}\right)
e^{i\int_{t_{i}}^{t_{f}}\mathcal{L}\left( \Psi\left( t^{\prime}\right) ,\dot{%
\Psi}\left( t^{\prime}\right) \right) dt^{\prime}} \\
& \left. \equiv\right. \left| \Phi\left( t_{f}\right) \right\rangle =\left|
\Psi\left( t_{f}\right) \right\rangle e^{i\mathcal{A}\left( \Psi,\tau\right)
}=\left( \frac{\left\langle \Psi\left( t_{i}\right) |\Psi\left( t_{i}\right)
\right\rangle }{\left\langle \Psi\left( t_{f}\right) |\Psi\left(
t_{f}\right) \right\rangle }\right) ^{\frac{1}{2}}\left| \Psi\left(
t_{f}\right) \right\rangle e^{i\func{Re}\mathcal{A}\left( \Psi,\tau\right) }
\end{align}
which we can see from 
\begin{align}
& i\int_{t_{i}}^{t_{f}}\mathcal{L}\left( \Psi\left( t^{\prime}\right) ,\dot{%
\Psi}\left( t^{\prime}\right) \right)
dt^{\prime}=i\int_{t_{i}}^{t_{f}}\left\{ \func{Re}\mathcal{L}\left(
\Psi\left( t^{\prime }\right) ,\dot{\Psi}\left( t^{\prime}\right) \right) +i%
\func{Im}\mathcal{L}\left( \Psi\left( t^{\prime}\right) ,\dot{\Psi}\left(
t^{\prime}\right) \right) \right\}  \notag \\
& =i\int_{t_{i}}^{t_{f}}\left\{ -\left\{ \frac{\func{Im}\left\langle
\Psi\left( t^{\prime}\right) |\dot{\Psi}\left( t^{\prime}\right)
\right\rangle }{\left\langle \Psi\left( t^{\prime}\right) |\Psi\left(
t^{\prime}\right) \right\rangle }+\frac{\left\langle \Psi\left( t^{\prime
}\right) |H\Psi\left( t^{\prime}\right) \right\rangle }{\left\langle
\Psi\left( t^{\prime}\right) |\Psi\left( t^{\prime}\right) \right\rangle }%
\right\} +\frac{i}{2}\frac{d\ln\left\langle \Psi\left( t^{\prime}\right)
|\Psi\left( t^{\prime}\right) \right\rangle }{dt^{\prime}}\right\}
dt^{\prime}  \notag \\
& =-i\int_{t_{i}}^{t_{f}}\left\{ \frac{\func{Im}\left\langle \Psi\left(
t^{\prime}\right) |\dot{\Psi}\left( t^{\prime}\right) \right\rangle }{%
\left\langle \Psi\left( t^{\prime}\right) |\Psi\left( t^{\prime}\right)
\right\rangle }+\frac{\left\langle \Psi\left( t^{\prime }\right)
|H\Psi\left( t^{\prime}\right) \right\rangle }{\left\langle \Psi\left(
t^{\prime}\right) |\Psi\left( t^{\prime}\right) \right\rangle }\right\}
dt^{\prime}-\frac{1}{2}\int_{t_{i}}^{t_{f}}\frac{d\ln\left\langle \Psi\left(
t^{\prime}\right) |\Psi\left( t^{\prime}\right) \right\rangle }{dt^{\prime}}%
dt^{\prime}  \notag \\
& =i\int_{t_{i}}^{t_{f}}\func{Re}\mathcal{L}\left( \Psi\left(
t^{\prime}\right) ,\dot{\Psi}\left( t^{\prime}\right) \right) dt^{\prime
}+\ln\left\{ \sqrt{\frac{\left\langle \Psi\left( t_{i}\right) |\Psi\left(
t_{i}\right) \right\rangle }{\left\langle \Psi\left( t_{f}\right)
|\Psi\left( t_{f}\right) \right\rangle }}\right\} =i\mathcal{A}\left(
\Psi,\tau\right)
\end{align}
thus 
\begin{equation}
\Phi\left( \mathsf{z,}t_{f}\right) =\left( \frac{\left\langle \Psi\left(
t_{i}\right) |\Psi\left( t_{i}\right) \right\rangle }{\left\langle
\Psi\left( t_{f}\right) |\Psi\left( t_{f}\right) \right\rangle }\right) ^{%
\frac{1}{2}}\Psi\left( \mathsf{z,}t_{f}\right) e^{i\func{Re}\mathcal{A}%
\left( \Psi,\tau\right) }  \label{phi}
\end{equation}
and in general 
\begin{equation}
\func{Re}\mathcal{A}\left( \Psi,\tau\right) \neq0
\end{equation}
We can see from eq. $\left( \ref{phi}\right) $ that the norm of the final
state is fixed by the choice of initial ''state'' but the phase is not. The
equation that determines the final phase is 
\begin{equation}
\frac{d\sigma}{dt}=-i\left\{ \frac{\func{Im}\left\langle \Psi\left( t\right)
|\dot{\Psi}\left( t\right) \right\rangle }{\left\langle \Psi\left( t\right)
|\Psi\left( t\right) \right\rangle }+\frac {\left\langle \Psi\left( t\right)
|H\Psi\left( t\right) \right\rangle }{\left\langle \Psi\left( t\right)
|\Psi\left( t\right) \right\rangle }\right\} =\func{Re}\mathcal{L}\left(
\Psi\left( t\right) ,\dot{\Psi}\left( t\right) \right) ;\sigma\left(
t_{i}\right) =0  \label{phase}
\end{equation}
where $\left| \Psi\left( t\right) \right\rangle $ is a solution of eq. $%
\left( \ref{preSE}\right) $ and the solution of eq. $\left( \ref{phase}%
\right) $ is 
\begin{equation}
\sigma\left( t\right) =\func{Re}\mathcal{A}\left( \Psi,\left[ t_{i},t\right]
\right)
\end{equation}

\subsection{Complex Scaling of the Action}

Transformations of the action 
\begin{equation}
\mathcal{A\longrightarrow A}_{F}
\end{equation}
and the substitution carried out to obtain the TDSE\ are special cases of
complex scaling of the trajectory, which we now consider in general. The
action is \emph{not invariant }to complex scaling (and therefore \emph{not}
a state function) and changes in the following manner%
\begin{align}
\mathcal{A}\left( \eta\Psi,\tau\right) & =\int_{t_{i}}^{t_{f}}\mathcal{\dot{A%
}}\left( \eta\left( t\right) \Psi\left( t\right) \right) dt  \notag \\
& =\int_{t_{i}}^{t_{f}}\left\{ \mathcal{\dot{A}}\left( \Psi\left( t\right)
\right) +i\frac{d\ln\eta\left( t\right) }{dt}\right\} dt=\mathcal{A}\left(
\Psi,\tau\right) +i\ln\left\{ \frac{\eta\left( t_{f}\right) }{\eta\left(
t_{i}\right) }\right\}
\end{align}
This property is deduced by first recalling 
\begin{equation}
\func{Re}\mathcal{L}\left( \Psi\left( t\right) ,\dot{\Psi}\left( t\right)
\right) =\frac{i}{2}\frac{\left\{ \left\langle \Psi\left( t\right) |\dot{\Psi%
}\left( t\right) \right\rangle -\left\langle \dot{\Psi }\left( t\right)
|\Psi\left( t\right) \right\rangle \right\} }{\left\langle \Psi\left(
t\right) |\Psi\left( t\right) \right\rangle }-\frac{\left\langle \Psi\left(
t\right) |H\Psi\left( t\right) \right\rangle }{\left\langle \Psi\left(
t\right) |\Psi\left( t\right) \right\rangle }
\end{equation}
which under the transformation of the paths 
\begin{equation}
\Phi\left( \mathsf{z,}t\right) =\eta\left( t\right) \Psi\left( \mathsf{z,}%
t\right)
\end{equation}
with 
\begin{equation}
\eta\left( t\right) =\left| \eta\left( t\right) \right| e^{i\sigma \left(
t\right) }=e^{\rho\left( t\right) }e^{i\sigma\left( t\right) }
\end{equation}
transforms as 
\begin{align}
& \func{Re}\mathcal{L}\left( \eta\left( t\right) \Psi\left( t\right) ,%
\overset{\bullet}{\overbrace{\eta\left( t\right) \Psi\left( t\right) }}%
\right)  \notag \\
& \left. =\right. \frac{i}{2}\frac{\left\{ \left\langle \Psi\left( t\right) |%
\frac{\partial}{\partial t}\left\{ \eta\left( t\right) \Psi\left( t\right)
\right\} \right\rangle -\left\langle \frac{\partial }{\partial t}\left\{
\eta\left( t\right) \Psi\left( t\right) \right\} |\Psi\left( t\right)
\right\rangle \right\} }{\left| \eta\left( t\right) \right| ^{2}\left\langle
\Psi\left( t\right) |\Psi\left( t\right) \right\rangle }-\frac{\left\langle
\Psi\left( t\right) |H\Psi\left( t\right) \right\rangle }{\left\langle
\Psi\left( t\right) |\Psi\left( t\right) \right\rangle }  \notag \\
& \left. {}\right. \qquad\qquad\qquad\lceil\frac{\partial}{\partial t}%
\left\{ \eta\left( t\right) \Psi\left( t\right) \right\} =\dot{\eta }\left(
t\right) \Psi\left( t\right) +\eta\left( t\right) \dot{\Psi }\left( t\right)
\rfloor  \notag \\
& \left. =\right. \frac{i}{2\left| \eta\left( t\right) \right|
^{2}\left\langle \Psi\left( t\right) |\Psi\left( t\right) \right\rangle }%
\left\{ \left\langle \Psi\left( t\right) |\Psi\left( t\right) \right\rangle
\eta\left( t\right) ^{\ast}\dot{\eta}\left( t\right) +\left\langle
\Psi\left( t\right) |\dot{\Psi}\left( t\right) \right\rangle \left|
\eta\left( t\right) \right| ^{2}\right.  \notag \\
& \qquad\qquad\qquad\left. -\left\langle \Psi\left( t\right) |\Psi\left(
t\right) \right\rangle \dot{\eta}\left( t\right) ^{\ast}\eta\left( t\right)
-\left\langle \dot{\Psi}\left( t\right) |\Psi\left( t\right) \right\rangle
\left| \eta\left( t\right) \right| ^{2}\right\} -\frac{\left\langle
\Psi\left( t\right) |H\Psi\left( t\right) \right\rangle }{\left\langle
\Psi\left( t\right) |\Psi\left( t\right) \right\rangle }  \notag \\
& \qquad\qquad\qquad\left. =\right. \func{Re}\mathcal{L}\left( \Psi\left(
t\right) ,\dot{\Psi}\left( t\right) \right) -\frac{i}{2\left| \eta\left(
t\right) \right| ^{2}}\left\{ \eta\left( t\right) ^{\ast}\dot{\eta}\left(
t\right) -\dot{\eta}\left( t\right) ^{\ast}\eta\left( t\right) \right\} 
\notag \\
& \left. {}\right. \lceil\frac{i}{2\left| \eta\left( t\right) \right| ^{2}}%
\left\{ \eta\left( t\right) ^{\ast}\dot{\eta}\left( t\right) -\dot{\eta}%
\left( t\right) ^{\ast}\eta\left( t\right) \right\} =\frac {i}{2}\left\{ 
\frac{\eta\left( t\right) ^{\ast}\dot{\eta}\left( t\right) }{\eta\left(
t\right) ^{\ast}\eta\left( t\right) }-\frac{\dot{\eta}\left( t\right)
^{\ast}\eta\left( t\right) }{\eta\left( t\right) ^{\ast}\eta\left( t\right) }%
\right\}  \notag \\
& \left. =\right. \frac{i}{2}\left\{ \frac{\dot{\eta}\left( t\right) }{%
\eta\left( t\right) }-\frac{\dot{\eta}\left( t\right) ^{\ast}}{\eta\left(
t\right) ^{\ast}}\right\} =\frac{i}{2}\frac{d}{dt}\left\{ \ln\eta\left(
t\right) -\ln\eta\left( t\right) ^{\ast}\right\} =\frac {i}{2}\frac{d}{dt}%
\ln\left\{ \frac{\eta\left( t\right) }{\eta\left( t\right) ^{\ast}}\right\} 
\notag \\
& =\frac{i}{2}\frac{d}{dt}\ln\left\{ \frac{e^{\rho\left( t\right)
}e^{i\sigma\left( t\right) }}{e^{\rho\left( t\right) }e^{-i\sigma\left(
t\right) }}\right\} =\frac{i}{2}\frac{d}{dt}\ln\left\{ \frac{e^{i\sigma
\left( t\right) }}{e^{-i\sigma\left( t\right) }}\right\} =\frac{i}{2}\frac{d%
}{dt}\ln e^{2i\sigma\left( t\right) }=-\frac{d\sigma\left( t\right) }{dt}%
\rfloor  \notag \\
& \Rightarrow\func{Re}\mathcal{L}\left( \eta\left( t\right) \Psi\left(
t\right) ,\overset{\bullet}{\overbrace{\eta\left( t\right) \Psi\left(
t\right) }}\right) =\func{Re}\mathcal{L}\left( \Psi\left( t\right) ,\dot{\Psi%
}\left( t\right) \right) -\frac {d\sigma\left( t\right) }{dt}
\end{align}
Now 
\begin{equation}
\func{Im}\mathcal{L}\left( \Psi\left( t\right) ,\dot{\Psi}\left( t\right)
\right) =\frac{i}{2}\frac{d\ln\left\langle \Psi\left( t\right) |\Psi\left(
t\right) \right\rangle }{dt}
\end{equation}
and 
\begin{align}
& \frac{d\ln\left\langle \eta\left( t\right) \Psi\left( t\right) |\eta\left(
t\right) \Psi\left( t\right) \right\rangle }{dt}=\frac {1}{\left\langle
\eta\left( t\right) \Psi\left( t\right) |\eta\left( t\right) \Psi\left(
t\right) \right\rangle }\frac{d\left\langle \eta\left( t\right) \Psi\left(
t\right) |\eta\left( t\right) \Psi\left( t\right) \right\rangle }{dt}  \notag
\\
& =\frac{1}{\left| \eta\left( t\right) \right| ^{2}\left\langle \Psi\left(
t\right) |\Psi\left( t\right) \right\rangle }\left\{ \frac{\left| \eta\left(
t\right) \right| ^{2}d\left\langle \Psi\left( t\right) |\Psi\left( t\right)
\right\rangle }{dt}+\left\langle \Psi\left( t\right) |\Psi\left( t\right)
\right\rangle \left\{ \dot{\eta}\left( t\right) ^{\ast}\eta\left( t\right)
+\eta\left( t\right) ^{\ast}\dot {\eta}\left( t\right) \right\} \right\} 
\notag \\
& \qquad\qquad\qquad\qquad\qquad\qquad\qquad\qquad\qquad\qquad\qquad\left.
=\right. \frac{d\ln\left\langle \Psi\left( t\right) |\Psi\left( t\right)
\right\rangle }{dt}+\frac{d\ln\left| \eta\left( t\right) \right| ^{2}}{dt}
\end{align}
hence%
\begin{align}
& \mathcal{L}\left( \eta\left( t\right) \Psi\left( t\right) ,\overset{\bullet%
}{\overbrace{\eta\left( t\right) \Psi\left( t\right) }}\right) =\func{Re}%
\mathcal{L}\left( \eta\left( t\right) \Psi\left( t\right) ,\overset{\bullet}{%
\overbrace{\eta\left( t\right) \Psi\left( t\right) }}\right) +i\func{Im}%
\mathcal{L}\left( \eta\left( t\right) \Psi\left( t\right) ,\overset{\bullet}{%
\overbrace {\eta\left( t\right) \Psi\left( t\right) }}\right)  \notag \\
& \qquad\qquad\left. =\right. \func{Re}\mathcal{L}\left( \Psi\left( t\right)
,\dot{\Psi}\left( t\right) \right) -\frac {d\sigma\left( t\right) }{dt}+%
\frac{i}{2}\left\{ \frac{d\ln\left\langle \Psi\left( t\right) |\Psi\left(
t\right) \right\rangle }{dt}+\frac {d\ln\left| \eta\left( t\right) \right|
^{2}}{dt}\right\}  \notag \\
& \qquad\qquad\qquad\qquad\qquad\qquad\qquad\qquad\qquad\left. =\right. 
\mathcal{L}\left( \Psi\left( t\right) ,\dot{\Psi}\left( t\right) \right) -%
\frac{d\sigma\left( t\right) }{dt}+i\frac{d\ln\left| \eta\left( t\right)
\right| }{dt}
\end{align}
If we let 
\begin{equation}
\left| \eta\left( t\right) \right| =e^{\rho\left( t\right) }
\end{equation}
one gets the more symmetric transformation%
\begin{align}
\mathcal{L}\left( \eta\left( t\right) \Psi\left( t\right) ,\overset{\bullet}{%
\overbrace{\eta\left( t\right) \Psi\left( t\right) }}\right) & =\mathcal{L}%
\left( \Psi\left( t\right) ,\dot{\Psi}\left( t\right) \right) -\frac{%
d\sigma\left( t\right) }{dt}+i\frac{d\rho\left( t\right) }{dt}  \notag \\
& =\mathcal{L}\left( \Psi\left( t\right) ,\dot{\Psi}\left( t\right) \right)
+i\frac{d\ln\eta\left( t\right) }{dt}  \label{Ltrans}
\end{align}
If $\eta\left( t\right) $ is fixed for $t\in\left( t_{i},t_{f}\right) $ then 
\begin{equation}
\frac{\delta}{\delta\eta}\mathcal{L}\left( \eta\left( t\right) \Psi\left(
t\right) ,\overset{\bullet}{\overbrace{\eta\left( t\right) \Psi\left(
t\right) }}\right) =0\quad\forall t\in\left( t_{i},t_{f}\right)
\end{equation}
and 
\begin{equation}
\frac{\delta}{\delta\Psi}\mathcal{L}\left( \eta\left( t\right) \Psi\left(
t\right) ,\overset{\bullet}{\overbrace{\eta\left( t\right) \Psi\left(
t\right) }}\right) =\frac{\delta}{\delta\Psi}\mathcal{L}\left( \Psi\left(
t\right) ,\dot{\Psi}\left( t\right) \right) \quad\forall t\in\left(
t_{i},t_{f}\right)  \label{inv}
\end{equation}
This transformation of the Lagrangian induces the following change in the
action%
\begin{equation}
\mathcal{A}\left( \eta\Psi,\tau\right) =\mathcal{A}\left( \Psi,\tau\right)
+i\ln\left\{ \frac{\eta\left( t_{f}\right) }{\eta\left( t_{i}\right) }%
\right\}
\end{equation}
and we note that $\eta\left( t_{i}\right) $ and $\eta\left( t_{f}\right) $
are complex variables and \emph{not} constant.

\subsection{Variable End Points and the Hamilton-Jacobi Equation}

The general principle of stationary action is obtained by setting the first
order variation 
\begin{align}
& \delta^{\left( 1\right) }\mathcal{A}\left( \Psi,\tau\right) =\int
_{t_{i}}^{t_{f}}\left\{ \frac{\delta\mathcal{L}\left( \Psi\left( t\right) ,%
\dot{\Psi}\left( t\right) \right) }{\delta\Psi\left( \mathsf{z}\right) }%
\delta\Psi\left( \mathsf{z}\right) +\frac{\delta\mathcal{L}\left( \Psi\left(
t\right) ,\dot{\Psi}\left( t\right) \right) }{\delta\dot{\Psi }\left( 
\mathsf{z}\right) }\delta\dot{\Psi}\left( \mathsf{z}\right) \right\} d%
\mathsf{z}dt  \notag \\
& =\int_{t_{i}}^{t_{f}}\left\{ \left\langle \left. \left[ \frac {\delta%
\mathcal{L}\left( \Psi\left( t\right) ,\dot{\Psi}\left( t\right) \right) }{%
\delta\Psi}\right] ^{\ast}\right| \delta\Psi\left( t\right) \right\rangle
+\left\langle \left. \left[ \frac{\delta\mathcal{L}\left( \Psi\left(
t\right) ,\dot{\Psi}\left( t\right) \right) }{\delta\dot{\Psi }}\right]
^{\ast}\right| \delta\dot{\Psi}\left( t\right) \right\rangle \right\} dt
\end{align}
equal to zero for all incremental changes, $\left| \delta\Psi\left( t\right)
\right\rangle $, of the trajectories $\left| \Psi\left( t\right)
\right\rangle $ \emph{and }changes of the time duration $\tau$ i.e. 
\begin{equation}
\delta^{\left( 1\right) }\mathcal{A}\left( \Psi,\tau\right) =\delta_{\Psi
}^{\left( 1\right) }\mathcal{A}\left( \Psi,\tau\right) +\delta_{\tau
}^{\left( 1\right) }\mathcal{A}\left( \Psi,\tau\right) =0
\end{equation}
As the initial and final state density operators are invariant to complex
scaling of the state vectors the initial and final state vectors are only
constrained to lie on the lines 
\begin{equation}
\left| \Psi\left( t_{\kappa}\right) \right\rangle =\eta\left( t_{\kappa
}\right) \left| \Lambda_{\kappa}\right\rangle ;\qquad\kappa=i,f;\quad
\eta\left( t_{\kappa}\right) \in\mathbb{C}  \label{const}
\end{equation}
where $\left| \Lambda_{\kappa}\right\rangle $ are fixed chosen end points
for $\kappa=i,f$. $\lceil$The amplitudes $\Lambda_{\kappa}\left( \mathsf{z}%
\right) =\left\langle \mathsf{z}|\Lambda_{\kappa}\right\rangle $ are fixed
functions of $\mathsf{z}$ which are in general different from each other$%
\rfloor$. For $t\in\left( t_{i},t_{f}\right) $ there are no constraints on
the state vectors. The scalar trajectories $\eta\left( t\right) \in\mathbb{C}
$ are arbitrary twice differentiable complex valued functions $\forall t\in%
\left[ t_{i},t_{f}\right] $. We have two families of spaces of $\mathcal{H}%
^{N}- $ valued incremental trajectories 
\begin{align}
\mathcal{S}_{0}\left( t_{f}\right) & =\left\{ \left| \delta\Psi\left(
t\right) \right\rangle \right\} ;\quad\left| \delta\Psi\left( t\right)
\right\rangle =0,\quad\kappa=i,f;\quad t\in\left[ t_{i},t_{f}\right] \\
\mathcal{S}_{1}\left( t_{f}\right) & =\left\{ \left| \delta\Psi\left(
t\right) \right\rangle \right\} ;\quad\left| \delta\Psi\left( t_{\kappa
}\right) \right\rangle =\eta\left( t_{\kappa}\right) \left| \Lambda
_{\kappa}\right\rangle ,\,\eta\left( t\right) \in\mathbb{C},\,\kappa
=i,f;\quad t\in\left[ t_{i},t_{f}\right]
\end{align}
clearly $\mathcal{S}_{1}\left( t_{f}\right) \supset\mathcal{S}_{0}\left(
t_{f}\right) $ as $\eta\left( t_{\kappa}\right) =0$ is always a possible
choice for the complex scaling parameter. The Lagrangian becomes a function
of $\eta\left( t_{\kappa}\right) $ on these lines given by 
\begin{align}
& \mathcal{L}\left( \Psi\left( \eta,t_{\kappa}\right) ,\dot{\Psi}\left(
\eta,t_{\kappa}\right) \right) =\mathcal{L}\left( \eta\left( t_{\kappa
}\right) ,\dot{\eta}\left( t_{\kappa}\right) ,\Lambda_{\kappa}\right)  \notag
\\
& =\frac{1}{\left| \eta\left( t_{\kappa}\right) \right| ^{2}\left\langle
\Lambda_{\kappa}|\Lambda_{\kappa}\right\rangle }\left\{ i\eta\left(
t_{\kappa}\right) ^{\ast}\left\langle \Lambda_{\kappa}\left| \overset{\cdot }%
{\overbrace{\eta\left( t_{\kappa}\right) \Psi\left( t_{\kappa}\right) }}%
\right. \right\rangle -\left| \eta\left( t_{\kappa}\right) \right|
^{2}\left\langle \Lambda_{\kappa}|H\Lambda_{\kappa}\right\rangle \right\} 
\notag \\
& =\frac{1}{\left| \eta\left( t_{\kappa}\right) \right| ^{2}\left\langle
\Lambda_{\kappa}|\Lambda_{\kappa}\right\rangle }\left\{ i\eta\left(
t_{\kappa}\right) ^{\ast}\left\langle \Lambda_{\kappa}\left| \left\{
\eta\left( t_{\kappa}\right) \dot{\Psi}\left( t_{\kappa}\right) +\dot {\eta}%
\left( t_{\kappa}\right) \Psi\left( t_{\kappa}\right) \right\} \right.
\right\rangle -\left| \eta\left( t_{\kappa}\right) \right| ^{2}\left\langle
\Lambda_{\kappa}|H\Lambda_{\kappa}\right\rangle \right\}  \notag \\
& =\frac{1}{\left| \eta\left( t_{\kappa}\right) \right| ^{2}\left\langle
\Lambda_{\kappa}|\Lambda_{\kappa}\right\rangle }\left\{ i\eta\left(
t_{\kappa}\right) ^{\ast}\left\{ \eta\left( t_{\kappa}\right) \left\langle
\Lambda_{\kappa}|\dot{\Psi}\left( t_{\kappa}\right) \right\rangle +\dot {\eta%
}\left( t_{\kappa}\right) \left\langle \Lambda_{\kappa}|\Psi\left(
t_{\kappa}\right) \right\rangle \right\} -\left| \eta\left( t_{\kappa
}\right) \right| ^{2}\left\langle \Lambda_{\kappa}|H\Lambda_{\kappa
}\right\rangle \right\}  \notag \\
& =\frac{1}{\left| \eta\left( t_{\kappa}\right) \right| ^{2}\left\langle
\Lambda_{\kappa}|\Lambda_{\kappa}\right\rangle }\left\{ i\eta\left(
t_{\kappa}\right) ^{\ast}\left\{ \eta\left( t_{\kappa}\right) \left\langle
\Lambda_{\kappa}|\dot{\Psi}\left( t_{\kappa}\right) \right\rangle +\dot {\eta%
}\left( t_{\kappa}\right) \eta\left( t_{\kappa}\right) \left\langle
\Lambda_{\kappa}|\Lambda_{\kappa}\right\rangle \right\} -\left| \eta\left(
t_{\kappa}\right) \right| ^{2}\left\langle \Lambda_{\kappa}|H\Lambda
_{\kappa}\right\rangle \right\}  \notag \\
& \qquad\qquad\qquad\qquad\qquad\qquad\qquad\qquad\qquad\left. =\right. 
\mathcal{L}\left( \Psi\left( t_{\kappa}\right) ,\dot{\Psi}\left(
t_{\kappa}\right) \right) +i\frac{d\ln\eta\left( t_{\kappa}\right) }{dt}
\end{align}
The stationarity condition for fixed a length of time process produces the
following equations for $\Psi\left( t\right) ,\,t\in\left[ t_{i},t_{f}\right]
$ 
\begin{align}
& \left\langle \left. \left[ \frac{\delta\mathcal{L}\left( \Psi\left(
t\right) ,\dot{\Psi}\left( t\right) \right) }{\delta\Psi}\right] ^{\ast
}\right| \delta\Psi\left( t\right) \right\rangle +\left\langle \left. \left[ 
\frac{\delta\mathcal{L}\left( \Psi\left( t\right) ,\dot{\Psi }\left(
t\right) \right) }{\delta\dot{\Psi}}\right] ^{\ast}\right| \delta\dot{\Psi}%
\left( t\right) \right\rangle =0  \notag \\
& \left. {}\right. \qquad\qquad\qquad\qquad\qquad\qquad\qquad\qquad
\qquad\qquad\forall\left| \delta\Psi\left( t\right) \right\rangle \in%
\mathcal{H}^{N}\,\forall t\in\left( t_{i},t_{f}\right)  \label{s1}
\end{align}
and more generally 
\begin{align}
& \int_{t_{i}}^{t_{f}}\left\{ \left\langle \left. \left[ \frac {\delta%
\mathcal{L}\left( \Psi\left( t\right) ,\dot{\Psi}\left( t\right) \right) }{%
\delta\Psi}\right] ^{\ast}\right| \delta\Psi\left( t\right) \right\rangle
+\left\langle \left. \left[ \frac{\delta\mathcal{L}\left( \Psi\left(
t\right) ,\dot{\Psi}\left( t\right) \right) }{\delta\dot{\Psi }}\right]
^{\ast}\right| \delta\dot{\Psi}\left( t\right) \right\rangle \right\} dt=0 
\notag \\
& \left. {}\right. \qquad\qquad\qquad\forall\left| \delta\Psi\left( t\right)
\right\rangle \in\mathcal{H}^{N}\,\forall t\in\left( t_{i},t_{f}\right)
;\qquad\left| \delta\Psi\left( t_{\kappa}\right) \right\rangle \in\mathcal{%
\eta}\left( t_{\kappa}\right) \left| \Lambda_{\kappa}\right\rangle
\,,\quad\kappa=i,f  \label{s2}
\end{align}
Eq. $\left( \ref{s1}\right) $ does not depend on $\left| \Psi\left(
t_{\kappa}\right) \right\rangle =$ $\left| \Lambda_{\kappa}\right\rangle $
(except as a fixed boundary condition) so it is sufficient to consider
incremental trajectories that belong to the subspace $\mathcal{S}_{0}$ that
have any values for $t\in\left( t_{i},t_{f}\right) $, i.e. $\left|
\delta\Psi\left( t\right) \right\rangle $ can be any vector belonging to $%
\mathcal{H}^{N}$ for $t\in\left( t_{i},t_{f}\right) $, and thus obtain the
usual Euler-Lagrange equations for $\left| \Psi\left( t\right) \right\rangle 
$ for the fixed end point conditions 
\begin{equation}
\left| \Psi\left( t_{\kappa}\right) \right\rangle =\left| \Lambda_{\kappa
}\right\rangle \quad\kappa=i,f.
\end{equation}
$\lceil$\textbf{Note }that these equations are $\emph{independent}$ of the
length of time duration $\tau$ i.e. the solutions only depend on the initial
and final points and the Hamiltonian of the process.$\rfloor$ Solutions to
the general problem must satisfy both eqs. $\left( \ref{s1}\right) $ and $%
\left( \ref{s2}\right) $ for variable time durations $\tau$. Noting that the
change of action induced by complex scaling of the initial and final point
is a function of the relative scaling of these two points one can regard the
action for the general problem as a function of the final position, $%
\eta\left( t_{f}\right) \left| \Lambda_{f}\right\rangle $, the time
interval, $\tau$, a \emph{fixed} initial point, $\Lambda_{\kappa}$, and of
stationary trajectories for the fixed end point problem.

The general equations are obtained by observing that the first order
variation $\delta_{\Psi}^{\left( 1\right) }\mathcal{A}$ of the action wrt
changes in stationary fixed point end point trajectories is 
\begin{align}
& \int_{t_{i}}^{t_{f}}\left\{ \left\langle \left. \left[ \frac {\delta%
\mathcal{L}\left( \Psi\left( t\right) ,\dot{\Psi}\left( t\right) \right) }{%
\delta\Psi}\right] ^{\ast}\right\vert \delta\Psi\left( t\right)
\right\rangle +\left\langle \left. \left[ \frac{\delta\mathcal{L}\left(
\Psi\left( t\right) ,\dot{\Psi}\left( t\right) \right) }{\delta\dot{\Psi }}%
\right] ^{\ast}\right\vert \delta\dot{\Psi}\left( t\right) \right\rangle
\right\} dt  \notag \\
& \left. =\right. \int_{t_{i}}^{t_{f}}\left\{ \left\langle \left. \left[ 
\frac{d}{dt}\left\{ \frac{\delta\mathcal{L}\left( \Psi\left( t\right) ,\dot{%
\Psi}\left( t\right) \right) }{\delta\dot{\Psi}}\right\} \right]
^{\ast}\right\vert \delta\Psi\left( t\right) \right\rangle \right.  \notag \\
& \qquad\qquad\qquad\qquad\qquad\qquad\qquad\qquad\qquad+\left. \left\langle
\left. \left[ \frac{\delta\mathcal{L}\left( \Psi\left( t\right) ,\dot{\Psi}%
\left( t\right) \right) }{\delta\dot{\Psi}}\right] ^{\ast }\right\vert \delta%
\dot{\Psi}\left( t\right) \right\rangle dt\right\}  \label{HJ1}
\end{align}
as the trajectories $\left\vert \Psi\left( t\right) \right\rangle $ satisfy
the Euler-Lagrange equation 
\begin{equation}
\frac{\delta\mathcal{L}\left( \Psi\left( t\right) ,\dot{\Psi}\left( t\right)
\right) }{\delta\Psi}=\frac{d}{dt}\left\{ \frac{\delta \mathcal{L}\left(
\Psi\left( t\right) ,\dot{\Psi}\left( t\right) \right) }{\delta\dot{\Psi}}%
\right\} \quad\forall t\in\left[ t_{i},t_{f}\right] .
\end{equation}
Noting that as $\left\vert \delta\Psi\left( t_{i}\right) \right\rangle =0$
the second term of eq. $\left( \ref{HJ1}\right) $ produces 
\begin{align}
& \int_{t_{i}}^{t_{f}}\left\langle \left. \left[ \frac{\delta\mathcal{L}%
\left( \Psi\left( t\right) ,\dot{\Psi}\left( t\right) \right) }{\delta\dot{%
\Psi}}\right] ^{\ast}\right\vert \delta\dot{\Psi}\left( t\right)
\right\rangle dt=\int_{t_{i}}^{t_{f}}\left\langle \left. \left[ \frac{\delta%
\mathcal{L}\left( \Psi\left( t\right) ,\dot{\Psi}\left( t\right) \right) }{%
\delta\dot{\Psi}}\right] ^{\ast}\right\vert \frac{d\delta\Psi\left( t\right) 
}{dt}\right\rangle dt  \notag \\
& =\left\langle \left. \left[ \frac{\delta\mathcal{L}\left( \Psi\left(
t_{f}\right) ,\dot{\Psi}\left( t_{f}\right) \right) }{\delta\dot{\Psi}}%
\right] ^{\ast}\right\vert \delta\Psi\left( t_{f}\right) \right\rangle
-\int_{t_{i}}^{t_{f}}\left\langle \left. \frac{d}{dt}\left\{ \frac {\delta%
\mathcal{L}\left( \Psi\left( t\right) ,\dot{\Psi}\left( t\right) \right) }{%
\delta\dot{\Psi}}\right\} ^{\ast}\right\vert \delta\Psi\left( t\right)
\right\rangle dt
\end{align}
Eq.$\left( \ref{HJ1}\right) $ then simplifies to 
\begin{equation}
\delta_{\Psi}^{\left( 1\right) }\mathcal{A}\left( \Psi,\tau\right)
=\left\langle \left. \left[ \frac{\delta\mathcal{L}\left( \Psi\left(
t_{f}\right) ,\dot{\Psi}\left( t_{f}\right) \right) }{\delta\dot{\Psi}}%
\right] ^{\ast}\right\vert \delta\Psi\left( t_{f}\right) \right\rangle
\end{equation}
The variations with respect to $t_{f}$ need to be done carefully. When we
fix the end point of the motion $\left\vert \Lambda_{f}\right\rangle $ but
change the arrival time $t_{f}$ to $t_{f}+\delta t_{f}$, we need to change $%
\left\vert \Psi\left( t_{f}\right) \right\rangle $ to $\left\vert \Psi\left(
t_{f}\right) \right\rangle -\left\vert \dot{\Psi}\left( t_{f}\right)
\right\rangle \delta t$ so that it arrives at $\left\vert
\Lambda_{f}\right\rangle $ at time $t_{f}+\delta t_{f}$ i.e. 
\begin{equation}
\delta_{\tau}\Psi\left( t_{f}\right) =-\left\vert \dot{\Psi}\left(
t_{f}\right) \right\rangle \delta t
\end{equation}
producing the change in the action 
\begin{equation}
-\left\langle \left. \left[ \frac{\delta\mathcal{L}\left( \Psi\left(
t_{f}\right) ,\dot{\Psi}\left( t_{f}\right) \right) }{\delta\dot{\Psi}}%
\right] ^{\ast}\right\vert \dot{\Psi}\left( t_{f}\right) \delta
t\right\rangle
\end{equation}
which must be added to the change induced by the change of length of process 
\begin{equation}
\mathcal{L}\left( \Psi\left( t_{f}\right) ,\dot{\Psi}\left( t_{f}\right)
\right) \delta t
\end{equation}
Adding these two together we obtain that the first order variation of the
action wrt time durations $\tau$ is 
\begin{equation}
\delta_{\tau}^{\left( 1\right) }\mathcal{A}\left( \Psi,\tau\right) =\mathcal{%
L}\left( \Psi\left( t_{f}\right) ,\dot{\Psi}\left( t_{f}\right) \right)
\delta t-\left\langle \left. \left[ \frac{\delta\mathcal{L}\left( \Psi\left(
t_{f}\right) ,\dot{\Psi}\left( t_{f}\right) \right) }{\delta\dot{\Psi}}%
\right] ^{\ast}\right\vert \dot{\Psi}\left( t_{f}\right) \delta
t\right\rangle
\end{equation}
The total first order change in the generalized stationary expression is
then given by 
\begin{align}
\delta^{\left( 1\right) }\mathcal{A}\left( \Psi,\tau\right) & =\left\langle
\left. \left[ \frac{\delta\mathcal{L}\left( \Psi\left( t_{f}\right) ,\dot{%
\Psi}\left( t_{f}\right) \right) }{\delta\dot{\Psi}}\right]
^{\ast}\right\vert \delta\Psi\left( t_{f}\right) \right\rangle +\mathcal{L}%
\left( \Psi\left( t_{f}\right) ,\dot{\Psi}\left( t_{f}\right) \right) \delta
t  \notag \\
& \qquad\qquad\qquad\qquad\qquad\qquad-\left\langle \left. \left[ \frac{%
\delta\mathcal{L}\left( \Psi\left( t_{f}\right) ,\dot{\Psi}\left(
t_{f}\right) \right) }{\delta\dot{\Psi}}\right] ^{\ast}\right\vert \dot{\Psi}%
\left( t_{f}\right) \delta t\right\rangle  \notag \\
& =\left\langle \left. \left[ \frac{\delta\mathcal{L}\left( \Psi\left(
t_{f}\right) ,\dot{\Psi}\left( t_{f}\right) \right) }{\delta\dot{\Psi}}%
\right] ^{\ast}\right\vert \delta\Psi\left( t_{f}\right) \right\rangle
-\left\langle \left. \left[ \frac{\delta\mathcal{L}\left( \Psi\left(
t_{f}\right) ,\dot{\Psi}\left( t_{f}\right) \right) }{\delta\dot{\Psi}}%
\right] ^{\ast}\right\vert \frac{d\Psi\left( t_{f}\right) }{dt}\delta
t\right\rangle  \notag \\
& \qquad\qquad\qquad\qquad\qquad\qquad\qquad\qquad\qquad\qquad\qquad \qquad+%
\mathcal{L}\left( \Psi\left( t_{f}\right) ,\dot{\Psi}\left( t_{f}\right)
\right) \delta t  \notag \\
& =\left\langle \left. \left[ \frac{\delta\mathcal{L}\left( \Psi\left(
t_{f}\right) ,\dot{\Psi}\left( t_{f}\right) \right) }{\delta\dot{\Psi}}%
\right] ^{\ast}\right\vert \delta\Psi\left( t_{f}\right) \right\rangle
-\left\langle \left. \left[ \frac{\delta\mathcal{L}\left( \Psi\left(
t_{f}\right) ,\dot{\Psi}\left( t_{f}\right) \right) }{\delta\dot{\Psi}}%
\right] ^{\ast}\right\vert \delta\Psi\left( t_{f}\right) \right\rangle 
\notag \\
& \qquad\qquad\qquad\qquad\qquad\qquad\qquad\qquad\qquad\qquad\qquad \qquad+%
\mathcal{L}\left( \Psi\left( t_{f}\right) ,\dot{\Psi}\left( t_{f}\right)
\right) \delta t
\end{align}
leading to 
\begin{equation}
\delta^{\left( 1\right) }\mathcal{A}\left( \Psi,\tau\right) =\mathcal{L}%
\left( \Psi\left( t_{f}\right) ,\dot{\Psi}\left( t_{f}\right) \right) \delta
t
\end{equation}
Now using the constraint $\left\vert \Psi\left( t_{f}\right) \right\rangle
=\eta\left( t_{f}\right) \left\vert \Lambda_{f}\right\rangle $ gives 
\begin{equation}
\delta^{\left( 1\right) }\mathcal{A}\left( \Psi,\tau\right) =\mathcal{L}%
\left( \eta\left( t_{f}\right) \left\vert \Lambda_{f}\right\rangle ,\left.
\left\vert \dot{\Psi}\left( t_{f}\right) \right\rangle \right\vert
_{\left\vert \Psi\left( t_{f}\right) \right\rangle =\eta\left( t_{f}\right)
\left\vert \Lambda_{f}\right\rangle }\right) +i\frac{d\ln \eta\left(
t_{\kappa}\right) }{dt}
\end{equation}
which on invoking stationary action produces the equation for the scaling
factor 
\begin{equation}
\frame{$\mathcal{L}\left( \left\vert \Lambda_f\right\rangle ,\left.
\left\vert \dot{\Psi}\left( t_f\right) \right\rangle \right\vert _\left\vert
\Psi\left( t_f\right) \right\rangle =\left\vert \Lambda_f\right\rangle
\right) +i\frac{d\ln\eta\left( t_\kappa\right) }{dt}$=0}
\end{equation}
that leads to the TDSE.

\end{document}
